% Options for packages loaded elsewhere
\PassOptionsToPackage{unicode}{hyperref}
\PassOptionsToPackage{hyphens}{url}
\documentclass[
]{article}
\usepackage{xcolor}
\usepackage[top=2cm,bottom=2cm,left=2cm,right=2cm]{geometry}
\usepackage{amsmath,amssymb}
\setcounter{secnumdepth}{-\maxdimen} % remove section numbering
\usepackage{iftex}
\ifPDFTeX
  \usepackage[T1]{fontenc}
  \usepackage[utf8]{inputenc}
  \usepackage{textcomp} % provide euro and other symbols
\else % if luatex or xetex
  \usepackage{unicode-math} % this also loads fontspec
  \defaultfontfeatures{Scale=MatchLowercase}
  \defaultfontfeatures[\rmfamily]{Ligatures=TeX,Scale=1}
\fi
\usepackage{lmodern}
\ifPDFTeX\else
  % xetex/luatex font selection
  \setmainfont[]{Times New Roman}
  \setmonofont[Scale=0.92]{Consolas}
\fi
% Use upquote if available, for straight quotes in verbatim environments
\IfFileExists{upquote.sty}{\usepackage{upquote}}{}
\IfFileExists{microtype.sty}{% use microtype if available
  \usepackage[]{microtype}
  \UseMicrotypeSet[protrusion]{basicmath} % disable protrusion for tt fonts
}{}
\makeatletter
\@ifundefined{KOMAClassName}{% if non-KOMA class
  \IfFileExists{parskip.sty}{%
    \usepackage{parskip}
  }{% else
    \setlength{\parindent}{0pt}
    \setlength{\parskip}{6pt plus 2pt minus 1pt}}
}{% if KOMA class
  \KOMAoptions{parskip=half}}
\makeatother
\usepackage{color}
\usepackage{fancyvrb}
\newcommand{\VerbBar}{|}
\newcommand{\VERB}{\Verb[commandchars=\\\{\}]}
\DefineVerbatimEnvironment{Highlighting}{Verbatim}{commandchars=\\\{\}}
% Add ',fontsize=\small' for more characters per line
\usepackage{framed}
\definecolor{shadecolor}{RGB}{248,248,248}
\newenvironment{Shaded}{\begin{snugshade}}{\end{snugshade}}
\newcommand{\AlertTok}[1]{\textcolor[rgb]{0.94,0.16,0.16}{#1}}
\newcommand{\AnnotationTok}[1]{\textcolor[rgb]{0.56,0.35,0.01}{\textbf{\textit{#1}}}}
\newcommand{\AttributeTok}[1]{\textcolor[rgb]{0.13,0.29,0.53}{#1}}
\newcommand{\BaseNTok}[1]{\textcolor[rgb]{0.00,0.00,0.81}{#1}}
\newcommand{\BuiltInTok}[1]{#1}
\newcommand{\CharTok}[1]{\textcolor[rgb]{0.31,0.60,0.02}{#1}}
\newcommand{\CommentTok}[1]{\textcolor[rgb]{0.56,0.35,0.01}{\textit{#1}}}
\newcommand{\CommentVarTok}[1]{\textcolor[rgb]{0.56,0.35,0.01}{\textbf{\textit{#1}}}}
\newcommand{\ConstantTok}[1]{\textcolor[rgb]{0.56,0.35,0.01}{#1}}
\newcommand{\ControlFlowTok}[1]{\textcolor[rgb]{0.13,0.29,0.53}{\textbf{#1}}}
\newcommand{\DataTypeTok}[1]{\textcolor[rgb]{0.13,0.29,0.53}{#1}}
\newcommand{\DecValTok}[1]{\textcolor[rgb]{0.00,0.00,0.81}{#1}}
\newcommand{\DocumentationTok}[1]{\textcolor[rgb]{0.56,0.35,0.01}{\textbf{\textit{#1}}}}
\newcommand{\ErrorTok}[1]{\textcolor[rgb]{0.64,0.00,0.00}{\textbf{#1}}}
\newcommand{\ExtensionTok}[1]{#1}
\newcommand{\FloatTok}[1]{\textcolor[rgb]{0.00,0.00,0.81}{#1}}
\newcommand{\FunctionTok}[1]{\textcolor[rgb]{0.13,0.29,0.53}{\textbf{#1}}}
\newcommand{\ImportTok}[1]{#1}
\newcommand{\InformationTok}[1]{\textcolor[rgb]{0.56,0.35,0.01}{\textbf{\textit{#1}}}}
\newcommand{\KeywordTok}[1]{\textcolor[rgb]{0.13,0.29,0.53}{\textbf{#1}}}
\newcommand{\NormalTok}[1]{#1}
\newcommand{\OperatorTok}[1]{\textcolor[rgb]{0.81,0.36,0.00}{\textbf{#1}}}
\newcommand{\OtherTok}[1]{\textcolor[rgb]{0.56,0.35,0.01}{#1}}
\newcommand{\PreprocessorTok}[1]{\textcolor[rgb]{0.56,0.35,0.01}{\textit{#1}}}
\newcommand{\RegionMarkerTok}[1]{#1}
\newcommand{\SpecialCharTok}[1]{\textcolor[rgb]{0.81,0.36,0.00}{\textbf{#1}}}
\newcommand{\SpecialStringTok}[1]{\textcolor[rgb]{0.31,0.60,0.02}{#1}}
\newcommand{\StringTok}[1]{\textcolor[rgb]{0.31,0.60,0.02}{#1}}
\newcommand{\VariableTok}[1]{\textcolor[rgb]{0.00,0.00,0.00}{#1}}
\newcommand{\VerbatimStringTok}[1]{\textcolor[rgb]{0.31,0.60,0.02}{#1}}
\newcommand{\WarningTok}[1]{\textcolor[rgb]{0.56,0.35,0.01}{\textbf{\textit{#1}}}}
\usepackage{longtable,booktabs,array}
\newcounter{none} % for unnumbered tables
\usepackage{calc} % for calculating minipage widths
% Correct order of tables after \paragraph or \subparagraph
\usepackage{etoolbox}
\makeatletter
\patchcmd\longtable{\par}{\if@noskipsec\mbox{}\fi\par}{}{}
\makeatother
% Allow footnotes in longtable head/foot
\IfFileExists{footnotehyper.sty}{\usepackage{footnotehyper}}{\usepackage{footnote}}
\makesavenoteenv{longtable}
\usepackage{graphicx}
\makeatletter
\newsavebox\pandoc@box
\newcommand*\pandocbounded[1]{% scales image to fit in text height/width
  \sbox\pandoc@box{#1}%
  \Gscale@div\@tempa{\textheight}{\dimexpr\ht\pandoc@box+\dp\pandoc@box\relax}%
  \Gscale@div\@tempb{\linewidth}{\wd\pandoc@box}%
  \ifdim\@tempb\p@<\@tempa\p@\let\@tempa\@tempb\fi% select the smaller of both
  \ifdim\@tempa\p@<\p@\scalebox{\@tempa}{\usebox\pandoc@box}%
  \else\usebox{\pandoc@box}%
  \fi%
}
% Set default figure placement to htbp
\def\fps@figure{htbp}
\makeatother
\setlength{\emergencystretch}{3em} % prevent overfull lines
\providecommand{\tightlist}{%
  \setlength{\itemsep}{0pt}\setlength{\parskip}{0pt}}
\usepackage{fancyhdr}
\usepackage{etoolbox}

% Fix longtable glue shrinkage in Pandoc tables
\AtBeginEnvironment{longtable}{%
  \setlength{\parskip}{0pt}%
}

% Page headers and footers
\pagestyle{fancy}
\fancyhf{}
\setlength{\headheight}{16pt}
\lhead{\footnotesize\textbf{DSAI1003 -- Fundamental Programming Concepts in Python}}
\rhead{\footnotesize\textbf{Lecture 2: Numbers, Strings and Operators}}
\lfoot{\scriptsize Faculty of Data Science and Artificial Intelligence -- National Economics University}
\rfoot{\footnotesize Page \thepage}
\renewcommand{\headrulewidth}{0.4pt}
\renewcommand{\footrulewidth}{0.3pt}
\usepackage{bookmark}
\IfFileExists{xurl.sty}{\usepackage{xurl}}{} % add URL line breaks if available
\urlstyle{same}
\hypersetup{
  pdftitle={Lesson 2: Programming with Numbers and Strings},
  pdfauthor={Faculty of Data Science \& Artificial Intelligence (FDA) - National Economics University},
  hidelinks,
  pdfcreator={LaTeX via pandoc}}

\title{Lesson 2: Programming with Numbers and Strings}
\author{Faculty of Data Science \& Artificial Intelligence (FDA) -
National Economics University}
\date{Course Code: DSAI1003}

\begin{document}
\maketitle

{
\setcounter{tocdepth}{3}
\tableofcontents
}
\section{Lesson 2. Programming with Numbers and
Strings}\label{lesson-2.-programming-with-numbers-and-strings}

\textbf{Last updated:} September 11, 2026

\begin{quote}
\textbf{Source:} Chapter 2 -- \emph{Programming with Numbers and
Strings}, in \emph{Python for Everyone, 3rd Edition} by Cay Horstmann
and Rance Necaise.

This lecture preserves the structure, terminology, examples, and
instructional focus of Chapter 2, while reorganizing the material for
direct classroom use. The examples intentionally follow the conventions
used in the textbook, including its variable-naming style and the
\texttt{\%} string-format operator.
\end{quote}

\begin{center}\rule{0.5\linewidth}{0.5pt}\end{center}

\subsection{Lesson Introduction}\label{lesson-introduction}

Chapter 1 introduced the idea of a computer program, the basic
organization of a computer, the Python programming environment, and the
role of algorithms. Chapter 2 begins the transition from
\emph{understanding what programs are} to \emph{writing programs that
perform useful computations}.

The main building blocks introduced in this chapter are:

\begin{enumerate}
\def\labelenumi{\arabic{enumi}.}
\tightlist
\item
  \textbf{Variables} for storing values.
\item
  \textbf{Numeric types} for representing integers and floating-point
  values.
\item
  \textbf{Arithmetic expressions} for computing new values.
\item
  \textbf{Strings} for representing and manipulating text.
\item
  \textbf{Input and output} for communicating with the program user.
\item
  \textbf{Simple graphics} for producing drawings with geometric shapes
  and text.
\end{enumerate}

A recurring theme throughout the chapter is that a program should not be
written by immediately typing Python statements. A better workflow is:

\begin{Shaded}
\begin{Highlighting}[]
\NormalTok{Understand the problem}
\NormalTok{        ↓}
\NormalTok{Work out a concrete example by hand}
\NormalTok{        ↓}
\NormalTok{Identify variables and constants}
\NormalTok{        ↓}
\NormalTok{Develop the computation / pseudocode}
\NormalTok{        ↓}
\NormalTok{Translate the solution into Python}
\NormalTok{        ↓}
\NormalTok{Run, inspect, and improve the result}
\end{Highlighting}
\end{Shaded}

\begin{center}\rule{0.5\linewidth}{0.5pt}\end{center}

\subsection{Knowledge and Skills to Be
Achieved}\label{knowledge-and-skills-to-be-achieved}

After completing this lesson, students will be able to:

\begin{itemize}
\tightlist
\item
  Define and use \textbf{variables} and \textbf{constants}.
\item
  Explain the difference between \texttt{int} and \texttt{float} values.
\item
  Apply the assignment operator \texttt{=} correctly.
\item
  Choose valid and descriptive variable names.
\item
  Use comments to explain program intent.
\item
  Write arithmetic expressions using \texttt{+}, \texttt{-}, \texttt{*},
  \texttt{/}, \texttt{//}, \texttt{\%}, and \texttt{**}.
\item
  Explain operator precedence and the role of parentheses.
\item
  Call built-in functions such as \texttt{abs()}, \texttt{round()},
  \texttt{min()}, and \texttt{max()}.
\item
  Import and use functions from the \texttt{math} module.
\item
  Recognize floating-point \textbf{roundoff errors}.
\item
  Perform a computation by hand before programming it.
\item
  Create and manipulate \textbf{strings}.
\item
  Use string concatenation, repetition, indexing, and methods.
\item
  Convert between strings and numbers using \texttt{str()},
  \texttt{int()}, and \texttt{float()}.
\item
  Use \texttt{input()} to obtain user input.
\item
  Format numeric and string output using the textbook's \texttt{\%}
  formatting operator.
\item
  Write small programs according to an \textbf{Input → Process → Output}
  structure.
\item
  Create simple graphical drawings using a graphics window and canvas.
\item
  Decompose a drawing into lines, rectangles, ovals, and text.
\end{itemize}

\begin{center}\rule{0.5\linewidth}{0.5pt}\end{center}

\subsection{Lesson Structure}\label{lesson-structure}

\begin{enumerate}
\def\labelenumi{\arabic{enumi}.}
\tightlist
\item
  Variables
\item
  Arithmetic
\item
  Problem Solving: First Do It By Hand
\item
  Strings
\item
  Input and Output
\item
  Graphics: Simple Drawings
\item
  Optional Toolbox: Symbolic Processing with SymPy
\item
  Summary and Exercises
\end{enumerate}

\begin{center}\rule{0.5\linewidth}{0.5pt}\end{center}

\section{2.1 Variables}\label{variables}

A program becomes useful when it can store values, modify them, and use
them in later computations.

A \textbf{variable} is a storage location in a computer program. Each
variable has:

\begin{itemize}
\tightlist
\item
  a \textbf{name};
\item
  a \textbf{value} currently stored in that location.
\end{itemize}

A useful analogy is a labeled parking space:

\pandocbounded{\includegraphics[keepaspectratio,alt={alt text}]{image.png}}
The name identifies the storage location; the value is its current
content.

\begin{center}\rule{0.5\linewidth}{0.5pt}\end{center}

\subsection{2.1.1 Defining Variables}\label{defining-variables}

A value is stored in a variable with an \textbf{assignment statement}.

\begin{Shaded}
\begin{Highlighting}[]
\NormalTok{cansPerPack }\OperatorTok{=} \DecValTok{6}
\end{Highlighting}
\end{Shaded}

The general form is:

\begin{Shaded}
\begin{Highlighting}[]
\NormalTok{variableName = value}
\end{Highlighting}
\end{Shaded}

The first time a variable is assigned a value, the variable is
\textbf{created and initialized}.

\begin{Shaded}
\begin{Highlighting}[]
\NormalTok{cansPerPack }\OperatorTok{=} \DecValTok{6}
\BuiltInTok{print}\NormalTok{(cansPerPack)}
\end{Highlighting}
\end{Shaded}

Output:

\begin{Shaded}
\begin{Highlighting}[]
\NormalTok{6}
\end{Highlighting}
\end{Shaded}

If the variable is assigned another value later, the previous value is
replaced.

\begin{Shaded}
\begin{Highlighting}[]
\NormalTok{cansPerPack }\OperatorTok{=} \DecValTok{6}
\NormalTok{cansPerPack }\OperatorTok{=} \DecValTok{8}
\BuiltInTok{print}\NormalTok{(cansPerPack)}
\end{Highlighting}
\end{Shaded}

Output:

\begin{Shaded}
\begin{Highlighting}[]
\NormalTok{8}
\end{Highlighting}
\end{Shaded}

\subsubsection{Assignment Is Not Mathematical
Equality}\label{assignment-is-not-mathematical-equality}

The symbol \texttt{=} in Python is the \textbf{assignment operator}. It
does not mean mathematical equality.

Consider:

\begin{Shaded}
\begin{Highlighting}[]
\NormalTok{cansPerPack }\OperatorTok{=}\NormalTok{ cansPerPack }\OperatorTok{+} \DecValTok{2}
\end{Highlighting}
\end{Shaded}

If \texttt{cansPerPack} currently contains \texttt{8}, Python performs
the following steps:

\begin{Shaded}
\begin{Highlighting}[]
\NormalTok{1. Read the current value: 8}
\NormalTok{2. Compute 8 + 2 = 10}
\NormalTok{3. Store 10 back into cansPerPack}
\end{Highlighting}
\end{Shaded}

After the statement executes:

\begin{Shaded}
\begin{Highlighting}[]
\NormalTok{cansPerPack }\OperatorTok{==} \DecValTok{10}
\end{Highlighting}
\end{Shaded}

conceptually, the variable now stores \texttt{10}.

This is perfectly meaningful in programming, even though the algebraic
statement \texttt{x\ =\ x\ +\ 2} would be impossible as an equation.

\begin{center}\rule{0.5\linewidth}{0.5pt}\end{center}

\subsection{2.1.2 Number Types}\label{number-types}

Every value in Python has a \textbf{data type}. The data type
determines:

\begin{itemize}
\tightlist
\item
  how the value is represented;
\item
  which operations can be performed on the value.
\end{itemize}

Two important numeric types are:

{\def\LTcaptype{none} % do not increment counter
\begin{longtable}[]{@{}lll@{}}
\toprule\noalign{}
Python type & Meaning & Examples \\
\midrule\noalign{}
\endhead
\bottomrule\noalign{}
\endlastfoot
\texttt{int} & integer / whole number & \texttt{6}, \texttt{0},
\texttt{-25} \\
\texttt{float} & floating-point number & \texttt{0.355}, \texttt{1.0},
\texttt{-3.5} \\
\end{longtable}
}

Examples:

\begin{Shaded}
\begin{Highlighting}[]
\NormalTok{count }\OperatorTok{=} \DecValTok{6}
\NormalTok{price }\OperatorTok{=} \FloatTok{2.95}
\end{Highlighting}
\end{Shaded}

A numeric value written directly in a program is called a \textbf{number
literal}.

\begin{Shaded}
\begin{Highlighting}[]
\DecValTok{6}
\FloatTok{0.5}
\FloatTok{1.0}
\FloatTok{1E6}
\FloatTok{2.96E{-}2}
\end{Highlighting}
\end{Shaded}

Important observations:

\begin{itemize}
\tightlist
\item
  \texttt{6} is an \texttt{int}.
\item
  \texttt{6.0} is a \texttt{float}.
\item
  Numbers written in exponential notation such as \texttt{1E6} are
  floating-point values.
\end{itemize}

The type is associated with the \textbf{value}, not permanently with the
variable name.

\begin{Shaded}
\begin{Highlighting}[]
\NormalTok{taxRate }\OperatorTok{=} \DecValTok{5}
\NormalTok{taxRate }\OperatorTok{=} \FloatTok{5.5}
\end{Highlighting}
\end{Shaded}

Python allows this. However, the chapter recommends keeping the type of
a variable conceptually consistent. If a variable represents a tax rate,
initialize it as a floating-point value when fractional values are
possible:

\begin{Shaded}
\begin{Highlighting}[]
\NormalTok{taxRate }\OperatorTok{=} \FloatTok{5.0}
\end{Highlighting}
\end{Shaded}

\begin{center}\rule{0.5\linewidth}{0.5pt}\end{center}

\subsection{2.1.3 Variable Names}\label{variable-names}

Variable names must follow Python's naming rules.

\subsubsection{Rules}\label{rules}

\begin{enumerate}
\def\labelenumi{\arabic{enumi}.}
\tightlist
\item
  A name must start with a letter or underscore \texttt{\_}.
\item
  Remaining characters may be letters, digits, or underscores.
\item
  Spaces and symbols such as \texttt{?}, \texttt{\%}, \texttt{.}, and
  \texttt{/} are not allowed.
\item
  Names are \textbf{case sensitive}.
\item
  Reserved words such as \texttt{if} and \texttt{class} cannot be used
  as variable names.
\end{enumerate}

Examples:

{\def\LTcaptype{none} % do not increment counter
\begin{longtable}[]{@{}
  >{\raggedright\arraybackslash}p{(\linewidth - 4\tabcolsep) * \real{0.3000}}
  >{\raggedleft\arraybackslash}p{(\linewidth - 4\tabcolsep) * \real{0.4000}}
  >{\raggedright\arraybackslash}p{(\linewidth - 4\tabcolsep) * \real{0.3000}}@{}}
\toprule\noalign{}
\begin{minipage}[b]{\linewidth}\raggedright
Name
\end{minipage} & \begin{minipage}[b]{\linewidth}\raggedleft
Valid?
\end{minipage} & \begin{minipage}[b]{\linewidth}\raggedright
Comment
\end{minipage} \\
\midrule\noalign{}
\endhead
\bottomrule\noalign{}
\endlastfoot
\texttt{canVolume1} & Yes & Letters and digits are allowed. \\
\texttt{x} & Yes & Legal, but often not descriptive enough. \\
\texttt{CanVolume} & Yes & Different from \texttt{canVolume}; violates
the book's lowercase-variable convention. \\
\texttt{6pack} & No & Cannot start with a digit. \\
\texttt{can\ volume} & No & Cannot contain spaces. \\
\texttt{class} & No & Reserved word. \\
\texttt{ltr/fl.oz} & No & Contains \texttt{/} and \texttt{.}. \\
\end{longtable}
}

\subsubsection{Descriptive Names}\label{descriptive-names}

Compare:

\begin{Shaded}
\begin{Highlighting}[]
\NormalTok{cv }\OperatorTok{=} \FloatTok{0.355}
\end{Highlighting}
\end{Shaded}

with:

\begin{Shaded}
\begin{Highlighting}[]
\NormalTok{canVolume }\OperatorTok{=} \FloatTok{0.355}
\end{Highlighting}
\end{Shaded}

The second name gives much more information about the meaning of the
variable.

\begin{center}\rule{0.5\linewidth}{0.5pt}\end{center}

\subsection{2.1.4 Constants}\label{constants}

A \textbf{constant} is a value that should remain unchanged throughout
the program.

Python does not technically prevent a programmer from changing a
constant. Instead, the chapter uses the convention of writing constant
names in uppercase letters.

\begin{Shaded}
\begin{Highlighting}[]
\NormalTok{BOTTLE\_VOLUME }\OperatorTok{=} \FloatTok{2.0}
\NormalTok{MAX\_SIZE }\OperatorTok{=} \DecValTok{100}
\end{Highlighting}
\end{Shaded}

Named constants make programs easier to understand.

Less clear:

\begin{Shaded}
\begin{Highlighting}[]
\NormalTok{totalVolume }\OperatorTok{=}\NormalTok{ bottles }\OperatorTok{*} \DecValTok{2}
\end{Highlighting}
\end{Shaded}

Clearer:

\begin{Shaded}
\begin{Highlighting}[]
\NormalTok{BOTTLE\_VOLUME }\OperatorTok{=} \FloatTok{2.0}
\NormalTok{totalVolume }\OperatorTok{=}\NormalTok{ bottles }\OperatorTok{*}\NormalTok{ BOTTLE\_VOLUME}
\end{Highlighting}
\end{Shaded}

The named constant explains \emph{why} the value \texttt{2.0} appears in
the computation.

\begin{center}\rule{0.5\linewidth}{0.5pt}\end{center}

\subsection{2.1.5 Comments}\label{comments}

A \textbf{comment} is an explanation written for human readers.

\begin{Shaded}
\begin{Highlighting}[]
\NormalTok{CAN\_VOLUME }\OperatorTok{=} \FloatTok{0.355}  \CommentTok{\# Liters in a 12{-}ounce can}
\end{Highlighting}
\end{Shaded}

Python ignores everything from \texttt{\#} to the end of the line.

A file-level comment can explain the purpose of a program:

\begin{Shaded}
\begin{Highlighting}[]
\CommentTok{\#\#}
\CommentTok{\# This program computes the volume, in liters, of a six{-}pack of soda cans.}
\CommentTok{\#}
\end{Highlighting}
\end{Shaded}

Comments are particularly useful for explaining:

\begin{itemize}
\tightlist
\item
  program purpose;
\item
  non-obvious constants;
\item
  important assumptions;
\item
  major steps of a computation.
\end{itemize}

\begin{center}\rule{0.5\linewidth}{0.5pt}\end{center}

\subsection{Example -- Computing Soda
Volume}\label{example-computing-soda-volume}

\begin{Shaded}
\begin{Highlighting}[]
\CommentTok{\#\#}
\CommentTok{\# This program computes the volume of a six{-}pack of soda cans}
\CommentTok{\# and the total volume when a two{-}liter bottle is added.}
\CommentTok{\#}

\NormalTok{CAN\_VOLUME }\OperatorTok{=} \FloatTok{0.355}
\NormalTok{BOTTLE\_VOLUME }\OperatorTok{=} \FloatTok{2.0}

\NormalTok{cansPerPack }\OperatorTok{=} \DecValTok{6}

\NormalTok{totalVolume }\OperatorTok{=}\NormalTok{ cansPerPack }\OperatorTok{*}\NormalTok{ CAN\_VOLUME}
\BuiltInTok{print}\NormalTok{(}\StringTok{"A six{-}pack of 12{-}ounce cans contains"}\NormalTok{, totalVolume, }\StringTok{"liters."}\NormalTok{)}

\NormalTok{totalVolume }\OperatorTok{=}\NormalTok{ totalVolume }\OperatorTok{+}\NormalTok{ BOTTLE\_VOLUME}
\BuiltInTok{print}\NormalTok{(}\StringTok{"A six{-}pack and a two{-}liter bottle contain"}\NormalTok{, totalVolume, }\StringTok{"liters."}\NormalTok{)}
\end{Highlighting}
\end{Shaded}

Conceptually:

\begin{Shaded}
\begin{Highlighting}[]
\NormalTok{6 cans × 0.355 L/can = 2.13 L}
\NormalTok{2.13 L + 2.00 L       = 4.13 L}
\end{Highlighting}
\end{Shaded}

\begin{center}\rule{0.5\linewidth}{0.5pt}\end{center}

\section{Common Error 2.1 -- Using Undefined
Variables}\label{common-error-2.1-using-undefined-variables}

A variable must be created before it is used.

Incorrect:

\begin{Shaded}
\begin{Highlighting}[]
\NormalTok{canVolume }\OperatorTok{=} \DecValTok{12} \OperatorTok{*}\NormalTok{ literPerOunce}
\NormalTok{literPerOunce }\OperatorTok{=} \FloatTok{0.0296}
\end{Highlighting}
\end{Shaded}

When Python executes the first statement, \texttt{literPerOunce} does
not yet exist.

Python therefore raises a \textbf{\texttt{NameError}}:

\begin{Shaded}
\begin{Highlighting}[]
\NormalTok{NameError: name \textquotesingle{}literPerOunce\textquotesingle{} is not defined}
\end{Highlighting}
\end{Shaded}

\subsection{Why Does This Error Occur?}\label{why-does-this-error-occur}

Python executes statements \textbf{from top to bottom}.

At this point:

\begin{Shaded}
\begin{Highlighting}[]
\NormalTok{canVolume }\OperatorTok{=} \DecValTok{12} \OperatorTok{*}\NormalTok{ literPerOunce}
\end{Highlighting}
\end{Shaded}

Python tries to find the current value of \texttt{literPerOunce}, but
that variable has not yet been assigned a value.

Therefore, the name is \textbf{undefined}.

Correct:

\begin{Shaded}
\begin{Highlighting}[]
\NormalTok{literPerOunce }\OperatorTok{=} \FloatTok{0.0296}
\NormalTok{canVolume }\OperatorTok{=} \DecValTok{12} \OperatorTok{*}\NormalTok{ literPerOunce}
\end{Highlighting}
\end{Shaded}

Program statements are executed in order, so definitions must appear
before uses.

\begin{center}\rule{0.5\linewidth}{0.5pt}\end{center}

\section{Programming Tip 2.1 -- Choose Descriptive Variable
Names}\label{programming-tip-2.1-choose-descriptive-variable-names}

Prefer:

\begin{Shaded}
\begin{Highlighting}[]
\NormalTok{canVolume }\OperatorTok{=} \FloatTok{0.355}
\end{Highlighting}
\end{Shaded}

over:

\begin{Shaded}
\begin{Highlighting}[]
\NormalTok{cv }\OperatorTok{=} \FloatTok{0.355}
\end{Highlighting}
\end{Shaded}

Descriptive names reduce the mental effort required to understand a
program and make code easier to maintain.

\begin{center}\rule{0.5\linewidth}{0.5pt}\end{center}

\section{Programming Tip 2.2 -- Do Not Use Magic
Numbers}\label{programming-tip-2.2-do-not-use-magic-numbers}

A \textbf{magic number} is an unexplained numeric literal in a program.

Less clear:

\begin{Shaded}
\begin{Highlighting}[]
\NormalTok{totalVolume }\OperatorTok{=}\NormalTok{ bottles }\OperatorTok{*} \DecValTok{2}
\end{Highlighting}
\end{Shaded}

Better:

\begin{Shaded}
\begin{Highlighting}[]
\NormalTok{BOTTLE\_VOLUME }\OperatorTok{=} \FloatTok{2.0}
\NormalTok{totalVolume }\OperatorTok{=}\NormalTok{ bottles }\OperatorTok{*}\NormalTok{ BOTTLE\_VOLUME}
\end{Highlighting}
\end{Shaded}

Using named constants makes changes safer. If the bottle size changes,
only the constant definition needs to be updated.

\begin{center}\rule{0.5\linewidth}{0.5pt}\end{center}

\subsection{Self Check 2.1}\label{self-check-2.1}

\subsubsection{Question 1}\label{question-1}

What happens the first time this statement is executed?

\begin{Shaded}
\begin{Highlighting}[]
\NormalTok{count }\OperatorTok{=} \DecValTok{10}
\end{Highlighting}
\end{Shaded}

Answer

A variable named \texttt{count} is created and initialized with the
integer value \texttt{10}.

\subsubsection{Question 2}\label{question-2}

Suppose:

\begin{Shaded}
\begin{Highlighting}[]
\NormalTok{x }\OperatorTok{=} \DecValTok{7}
\NormalTok{x }\OperatorTok{=}\NormalTok{ x }\OperatorTok{+} \DecValTok{3}
\end{Highlighting}
\end{Shaded}

What value does \texttt{x} contain afterwards?

Answer

\texttt{10}.

Python first evaluates \texttt{x\ +\ 3}, then stores the result back
into \texttt{x}.

\subsubsection{Question 3}\label{question-3}

What is the difference between \texttt{6} and \texttt{6.0}?

Answer

\texttt{6} has type \texttt{int}; \texttt{6.0} has type \texttt{float}.

\subsubsection{Question 4}\label{question-4}

Which name best represents a constant?

A. \texttt{bottleVolume}\\
B. \texttt{BOTTLE\_VOLUME}\\
C. \texttt{bottle\ volume}\\
D. \texttt{2liters}

Answer

\textbf{B. \texttt{BOTTLE\_VOLUME}}

The chapter uses uppercase names for constants.

\subsubsection{Question 5}\label{question-5}

Why is the following code incorrect?

\begin{Shaded}
\begin{Highlighting}[]
\NormalTok{result }\OperatorTok{=}\NormalTok{ price }\OperatorTok{*}\NormalTok{ quantity}
\NormalTok{quantity }\OperatorTok{=} \DecValTok{4}
\end{Highlighting}
\end{Shaded}

Answer

\texttt{quantity} is used before it is created and initialized.

\begin{center}\rule{0.5\linewidth}{0.5pt}\end{center}

\section{2.2 Arithmetic}\label{arithmetic}

Python can be used as a calculator, but arithmetic expressions must be
written using Python's operators and syntax.

\begin{center}\rule{0.5\linewidth}{0.5pt}\end{center}

\subsection{2.2.1 Basic Arithmetic
Operations}\label{basic-arithmetic-operations}

The basic operators are:

{\def\LTcaptype{none} % do not increment counter
\begin{longtable}[]{@{}lrl@{}}
\toprule\noalign{}
Operation & Python operator & Example \\
\midrule\noalign{}
\endhead
\bottomrule\noalign{}
\endlastfoot
Addition & \texttt{+} & \texttt{a\ +\ b} \\
Subtraction & \texttt{-} & \texttt{a\ -\ b} \\
Multiplication & \texttt{*} & \texttt{a\ *\ b} \\
Division & \texttt{/} & \texttt{a\ /\ b} \\
\end{longtable}
}

For example, the mathematical expression

\[
\frac{a+b}{2}
\]

is written in Python as:

\begin{Shaded}
\begin{Highlighting}[]
\NormalTok{(a }\OperatorTok{+}\NormalTok{ b) }\OperatorTok{/} \DecValTok{2}
\end{Highlighting}
\end{Shaded}

The combination of literals, variables, operators, and parentheses is
called an \textbf{expression}.

\subsubsection{Operator Precedence}\label{operator-precedence}

Python follows precedence rules similar to algebra:

\begin{enumerate}
\def\labelenumi{\arabic{enumi}.}
\tightlist
\item
  Parentheses
\item
  Exponentiation
\item
  Multiplication and division
\item
  Addition and subtraction
\end{enumerate}

Example:

\begin{Shaded}
\begin{Highlighting}[]
\NormalTok{a }\OperatorTok{+}\NormalTok{ b }\OperatorTok{/} \DecValTok{2}
\end{Highlighting}
\end{Shaded}

computes \texttt{b\ /\ 2} first.

To average \texttt{a} and \texttt{b}, use:

\begin{Shaded}
\begin{Highlighting}[]
\NormalTok{(a }\OperatorTok{+}\NormalTok{ b) }\OperatorTok{/} \DecValTok{2}
\end{Highlighting}
\end{Shaded}

Operators of the same precedence are generally evaluated from left to
right.

\begin{Shaded}
\begin{Highlighting}[]
\DecValTok{10} \OperatorTok{{-}} \DecValTok{2} \OperatorTok{{-}} \DecValTok{3}
\end{Highlighting}
\end{Shaded}

is evaluated as:

\begin{Shaded}
\begin{Highlighting}[]
\NormalTok{(10 {-} 2) {-} 3 = 5}
\end{Highlighting}
\end{Shaded}

\subsubsection{\texorpdfstring{Mixing \texttt{int} and
\texttt{float}}{Mixing int and float}}\label{mixing-int-and-float}

If integer and floating-point values are mixed in an arithmetic
expression, the result is a floating-point value.

\begin{Shaded}
\begin{Highlighting}[]
\DecValTok{7} \OperatorTok{+} \FloatTok{4.0}
\end{Highlighting}
\end{Shaded}

Result:

\begin{Shaded}
\begin{Highlighting}[]
\NormalTok{11.0}
\end{Highlighting}
\end{Shaded}

\begin{center}\rule{0.5\linewidth}{0.5pt}\end{center}

\subsection{2.2.2 Powers}\label{powers}

Python uses \texttt{**} for exponentiation.

\begin{Shaded}
\begin{Highlighting}[]
\DecValTok{2} \OperatorTok{**} \DecValTok{3}
\end{Highlighting}
\end{Shaded}

means:

\[
2^3 = 8
\]

Example:

\begin{Shaded}
\begin{Highlighting}[]
\DecValTok{10} \OperatorTok{*} \DecValTok{2} \OperatorTok{**} \DecValTok{3}
\end{Highlighting}
\end{Shaded}

Result:

\begin{Shaded}
\begin{Highlighting}[]
\NormalTok{80}
\end{Highlighting}
\end{Shaded}

Exponentiation has higher precedence than multiplication.

Power operations are evaluated from right to left:

\begin{Shaded}
\begin{Highlighting}[]
\DecValTok{10} \OperatorTok{**} \DecValTok{2} \OperatorTok{**} \DecValTok{3}
\end{Highlighting}
\end{Shaded}

is interpreted as:

\begin{Shaded}
\begin{Highlighting}[]
\NormalTok{10 ** (2 ** 3)}
\end{Highlighting}
\end{Shaded}

not:

\begin{Shaded}
\begin{Highlighting}[]
\NormalTok{(10 ** 2) ** 3}
\end{Highlighting}
\end{Shaded}

A mathematical expression such as

\[
b\left(1 + \frac{r}{100}\right)^n
\]

becomes:

\begin{Shaded}
\begin{Highlighting}[]
\NormalTok{b }\OperatorTok{*}\NormalTok{ (}\DecValTok{1} \OperatorTok{+}\NormalTok{ r }\OperatorTok{/} \DecValTok{100}\NormalTok{) }\OperatorTok{**}\NormalTok{ n}
\end{Highlighting}
\end{Shaded}

\begin{center}\rule{0.5\linewidth}{0.5pt}\end{center}

\subsection{2.2.3 Floor Division and
Remainder}\label{floor-division-and-remainder}

Normal division \texttt{/} produces a floating-point result.

\begin{Shaded}
\begin{Highlighting}[]
\DecValTok{7} \OperatorTok{/} \DecValTok{4}
\end{Highlighting}
\end{Shaded}

Result:

\begin{Shaded}
\begin{Highlighting}[]
\NormalTok{1.75}
\end{Highlighting}
\end{Shaded}

\subsubsection{\texorpdfstring{Floor Division
\texttt{//}}{Floor Division //}}\label{floor-division}

For positive integers, \texttt{//} computes the quotient and discards
the fractional part.

\begin{Shaded}
\begin{Highlighting}[]
\DecValTok{7} \OperatorTok{//} \DecValTok{4}
\end{Highlighting}
\end{Shaded}

Result:

\begin{Shaded}
\begin{Highlighting}[]
\NormalTok{1}
\end{Highlighting}
\end{Shaded}

\subsubsection{\texorpdfstring{Remainder
\texttt{\%}}{Remainder \%}}\label{remainder}

The \texttt{\%} operator computes the remainder.

\begin{Shaded}
\begin{Highlighting}[]
\DecValTok{7} \OperatorTok{\%} \DecValTok{4}
\end{Highlighting}
\end{Shaded}

Result:

\begin{Shaded}
\begin{Highlighting}[]
\NormalTok{3}
\end{Highlighting}
\end{Shaded}

\subsubsection{Example -- Converting Pennies to Dollars and
Cents}\label{example-converting-pennies-to-dollars-and-cents}

\begin{Shaded}
\begin{Highlighting}[]
\NormalTok{pennies }\OperatorTok{=} \DecValTok{1729}

\NormalTok{dollars }\OperatorTok{=}\NormalTok{ pennies }\OperatorTok{//} \DecValTok{100}
\NormalTok{cents }\OperatorTok{=}\NormalTok{ pennies }\OperatorTok{\%} \DecValTok{100}

\BuiltInTok{print}\NormalTok{(dollars)}
\BuiltInTok{print}\NormalTok{(cents)}
\end{Highlighting}
\end{Shaded}

Output:

\begin{Shaded}
\begin{Highlighting}[]
\NormalTok{17}
\NormalTok{29}
\end{Highlighting}
\end{Shaded}

Useful patterns for a positive integer \texttt{n}:

\begin{Shaded}
\begin{Highlighting}[]
\NormalTok{n }\OperatorTok{\%} \DecValTok{10}       \CommentTok{\# last digit}
\NormalTok{n }\OperatorTok{//} \DecValTok{10}      \CommentTok{\# all digits except the last}
\NormalTok{n }\OperatorTok{\%} \DecValTok{100}      \CommentTok{\# last two digits}
\NormalTok{n }\OperatorTok{\%} \DecValTok{2}        \CommentTok{\# 0 if even, 1 if odd}
\end{Highlighting}
\end{Shaded}

\begin{center}\rule{0.5\linewidth}{0.5pt}\end{center}

\subsection{2.2.4 Calling Functions}\label{calling-functions}

A function may \textbf{return a value} that can be used in an
expression, printed, or stored in a variable.

Example:

\begin{Shaded}
\begin{Highlighting}[]
\NormalTok{distance }\OperatorTok{=} \BuiltInTok{abs}\NormalTok{(}\OperatorTok{{-}}\DecValTok{173}\NormalTok{)}
\end{Highlighting}
\end{Shaded}

\texttt{abs()} returns the absolute value.

Other useful built-in mathematical functions include:

{\def\LTcaptype{none} % do not increment counter
\begin{longtable}[]{@{}ll@{}}
\toprule\noalign{}
Function & Purpose \\
\midrule\noalign{}
\endhead
\bottomrule\noalign{}
\endlastfoot
\texttt{abs(x)} & Absolute value of \texttt{x} \\
\texttt{round(x)} & Round \texttt{x} to a whole number \\
\texttt{round(x,\ n)} & Round \texttt{x} to \texttt{n} decimal places \\
\texttt{max(x1,\ ...,\ xn)} & Largest argument \\
\texttt{min(x1,\ ...,\ xn)} & Smallest argument \\
\end{longtable}
}

Examples:

\begin{Shaded}
\begin{Highlighting}[]
\BuiltInTok{print}\NormalTok{(}\BuiltInTok{abs}\NormalTok{(}\OperatorTok{{-}}\DecValTok{10}\NormalTok{))}
\BuiltInTok{print}\NormalTok{(}\BuiltInTok{round}\NormalTok{(}\FloatTok{7.627}\NormalTok{, }\DecValTok{2}\NormalTok{))}
\BuiltInTok{print}\NormalTok{(}\BuiltInTok{min}\NormalTok{(}\FloatTok{7.25}\NormalTok{, }\FloatTok{10.95}\NormalTok{, }\FloatTok{5.95}\NormalTok{, }\FloatTok{6.05}\NormalTok{))}
\end{Highlighting}
\end{Shaded}

A function must be called with an appropriate number of arguments.

Incorrect:

\begin{Shaded}
\begin{Highlighting}[]
\BuiltInTok{abs}\NormalTok{()}
\BuiltInTok{abs}\NormalTok{(}\OperatorTok{{-}}\DecValTok{10}\NormalTok{, }\DecValTok{2}\NormalTok{)}
\end{Highlighting}
\end{Shaded}

The \texttt{abs()} function expects exactly one argument.

\begin{center}\rule{0.5\linewidth}{0.5pt}\end{center}

\subsection{\texorpdfstring{2.2.5 Mathematical Functions and the
\texttt{math}
Module}{2.2.5 Mathematical Functions and the math Module}}\label{mathematical-functions-and-the-math-module}

Python's standard library contains reusable modules.

The \texttt{math} module provides many mathematical functions.

To use \texttt{sqrt()}:

\begin{Shaded}
\begin{Highlighting}[]
\ImportTok{from}\NormalTok{ math }\ImportTok{import}\NormalTok{ sqrt}

\NormalTok{y }\OperatorTok{=}\NormalTok{ sqrt(}\DecValTok{25}\NormalTok{)}
\BuiltInTok{print}\NormalTok{(y)}
\end{Highlighting}
\end{Shaded}

Output:

\begin{Shaded}
\begin{Highlighting}[]
\NormalTok{5.0}
\end{Highlighting}
\end{Shaded}

Selected functions include:

{\def\LTcaptype{none} % do not increment counter
\begin{longtable}[]{@{}ll@{}}
\toprule\noalign{}
Function & Meaning \\
\midrule\noalign{}
\endhead
\bottomrule\noalign{}
\endlastfoot
\texttt{sqrt(x)} & Square root \\
\texttt{trunc(x)} & Truncate floating-point value to an integer \\
\texttt{sin(x)} & Sine of \texttt{x} in radians \\
\texttt{cos(x)} & Cosine of \texttt{x} in radians \\
\texttt{tan(x)} & Tangent of \texttt{x} in radians \\
\texttt{exp(x)} & \(e^x\) \\
\texttt{degrees(x)} & Radians → degrees \\
\texttt{radians(x)} & Degrees → radians \\
\texttt{log(x)} & Natural logarithm \\
\texttt{log(x,\ base)} & Logarithm with specified base \\
\end{longtable}
}

The module also provides constants such as \texttt{pi}.

\begin{Shaded}
\begin{Highlighting}[]
\ImportTok{from}\NormalTok{ math }\ImportTok{import}\NormalTok{ pi}

\NormalTok{area }\OperatorTok{=}\NormalTok{ pi }\OperatorTok{*}\NormalTok{ radius }\OperatorTok{**} \DecValTok{2}
\end{Highlighting}
\end{Shaded}

\begin{center}\rule{0.5\linewidth}{0.5pt}\end{center}

\section{Common Error 2.2 -- Roundoff
Errors}\label{common-error-2.2-roundoff-errors}

Floating-point numbers cannot always represent decimal values exactly.

Example:

\begin{Shaded}
\begin{Highlighting}[]
\NormalTok{price }\OperatorTok{=} \FloatTok{4.35}
\NormalTok{quantity }\OperatorTok{=} \DecValTok{100}
\NormalTok{total }\OperatorTok{=}\NormalTok{ price }\OperatorTok{*}\NormalTok{ quantity}
\BuiltInTok{print}\NormalTok{(total)}
\end{Highlighting}
\end{Shaded}

A computer may display a value close to, but not exactly, \texttt{435}.

This is not necessarily a programming mistake. It is a consequence of
representing many decimal fractions using a finite number of binary
digits.

For user-facing output, round the value or format it with a fixed number
of decimal places.

\begin{center}\rule{0.5\linewidth}{0.5pt}\end{center}

\section{Common Error 2.3 -- Unbalanced
Parentheses}\label{common-error-2.3-unbalanced-parentheses}

Consider:

\begin{Shaded}
\begin{Highlighting}[]
\NormalTok{((a }\OperatorTok{+}\NormalTok{ b) }\OperatorTok{*}\NormalTok{ t }\OperatorTok{/} \DecValTok{2} \OperatorTok{*}\NormalTok{ (}\DecValTok{1} \OperatorTok{{-}}\NormalTok{ t)}
\end{Highlighting}
\end{Shaded}

There are more opening parentheses than closing parentheses.

A useful checking rule is:

\begin{itemize}
\tightlist
\item
  while scanning left to right, the number of closing parentheses should
  never exceed the number of opening parentheses;
\item
  at the end, the counts must match.
\end{itemize}

\begin{center}\rule{0.5\linewidth}{0.5pt}\end{center}

\section{Programming Tip 2.3 -- Use Spaces in
Expressions}\label{programming-tip-2.3-use-spaces-in-expressions}

Prefer:

\begin{Shaded}
\begin{Highlighting}[]
\NormalTok{x1 }\OperatorTok{=}\NormalTok{ (}\OperatorTok{{-}}\NormalTok{b }\OperatorTok{+}\NormalTok{ sqrt(b }\OperatorTok{**} \DecValTok{2} \OperatorTok{{-}} \DecValTok{4} \OperatorTok{*}\NormalTok{ a }\OperatorTok{*}\NormalTok{ c)) }\OperatorTok{/}\NormalTok{ (}\DecValTok{2} \OperatorTok{*}\NormalTok{ a)}
\end{Highlighting}
\end{Shaded}

over:

\begin{Shaded}
\begin{Highlighting}[]
\NormalTok{x1}\OperatorTok{=}\NormalTok{(}\OperatorTok{{-}}\NormalTok{b}\OperatorTok{+}\NormalTok{sqrt(b}\OperatorTok{**}\DecValTok{2}\OperatorTok{{-}}\DecValTok{4}\OperatorTok{*}\NormalTok{a}\OperatorTok{*}\NormalTok{c))}\OperatorTok{/}\NormalTok{(}\DecValTok{2}\OperatorTok{*}\NormalTok{a)}
\end{Highlighting}
\end{Shaded}

Spaces make the structure of an expression easier to read.

\begin{center}\rule{0.5\linewidth}{0.5pt}\end{center}

\section{Special Topic 2.1 -- Other Ways to Import
Modules}\label{special-topic-2.1-other-ways-to-import-modules}

Import selected functions:

\begin{Shaded}
\begin{Highlighting}[]
\ImportTok{from}\NormalTok{ math }\ImportTok{import}\NormalTok{ sqrt, sin, cos}
\end{Highlighting}
\end{Shaded}

Import all names from a module:

\begin{Shaded}
\begin{Highlighting}[]
\ImportTok{from}\NormalTok{ math }\ImportTok{import} \OperatorTok{*}
\end{Highlighting}
\end{Shaded}

Import the module itself:

\begin{Shaded}
\begin{Highlighting}[]
\ImportTok{import}\NormalTok{ math}

\NormalTok{y }\OperatorTok{=}\NormalTok{ math.sqrt(x)}
\end{Highlighting}
\end{Shaded}

The last style makes the origin of the function explicit.

\begin{center}\rule{0.5\linewidth}{0.5pt}\end{center}

\section{Special Topic 2.2 -- Combining Assignment and
Arithmetic}\label{special-topic-2.2-combining-assignment-and-arithmetic}

Python provides augmented assignment operators.

\begin{Shaded}
\begin{Highlighting}[]
\NormalTok{total }\OperatorTok{+=}\NormalTok{ cans}
\end{Highlighting}
\end{Shaded}

is equivalent to:

\begin{Shaded}
\begin{Highlighting}[]
\NormalTok{total }\OperatorTok{=}\NormalTok{ total }\OperatorTok{+}\NormalTok{ cans}
\end{Highlighting}
\end{Shaded}

Similarly:

\begin{Shaded}
\begin{Highlighting}[]
\NormalTok{total }\OperatorTok{*=} \DecValTok{2}
\NormalTok{count }\OperatorTok{+=} \DecValTok{1}
\end{Highlighting}
\end{Shaded}

\begin{center}\rule{0.5\linewidth}{0.5pt}\end{center}

\section{Special Topic 2.3 -- Line
Joining}\label{special-topic-2.3-line-joining}

Long expressions can be split across lines when the break occurs inside
parentheses.

\begin{Shaded}
\begin{Highlighting}[]
\NormalTok{x1 }\OperatorTok{=}\NormalTok{ ((}\OperatorTok{{-}}\NormalTok{b }\OperatorTok{+}\NormalTok{ sqrt(b }\OperatorTok{**} \DecValTok{2} \OperatorTok{{-}} \DecValTok{4} \OperatorTok{*}\NormalTok{ a }\OperatorTok{*}\NormalTok{ c))}
      \OperatorTok{/}\NormalTok{ (}\DecValTok{2} \OperatorTok{*}\NormalTok{ a))}
\end{Highlighting}
\end{Shaded}

This is preferable to breaking a complete expression at an arbitrary
point.

\begin{center}\rule{0.5\linewidth}{0.5pt}\end{center}

\subsection{Self Check 2.2}\label{self-check-2.2}

\subsubsection{Question 1}\label{question-1-1}

What is the value of:

\begin{Shaded}
\begin{Highlighting}[]
\DecValTok{10} \OperatorTok{+} \DecValTok{6} \OperatorTok{/} \DecValTok{2}
\end{Highlighting}
\end{Shaded}

Answer

\texttt{13.0}

Division is performed before addition.

\subsubsection{Question 2}\label{question-2-1}

What is the result of:

\begin{Shaded}
\begin{Highlighting}[]
\DecValTok{17} \OperatorTok{//} \DecValTok{5}
\DecValTok{17} \OperatorTok{\%} \DecValTok{5}
\end{Highlighting}
\end{Shaded}

Answer

\begin{Shaded}
\begin{Highlighting}[]
\NormalTok{3}
\NormalTok{2}
\end{Highlighting}
\end{Shaded}

\subsubsection{Question 3}\label{question-3-1}

What is the last digit of a positive integer \texttt{n}?

Answer

\begin{Shaded}
\begin{Highlighting}[]
\NormalTok{n }\OperatorTok{\%} \DecValTok{10}
\end{Highlighting}
\end{Shaded}

\subsubsection{Question 4}\label{question-4-1}

Why is the following expression different from \texttt{(a\ +\ b)\ /\ 2}?

\begin{Shaded}
\begin{Highlighting}[]
\NormalTok{a }\OperatorTok{+}\NormalTok{ b }\OperatorTok{/} \DecValTok{2}
\end{Highlighting}
\end{Shaded}

Answer

Because division has higher precedence than addition. Python computes
\texttt{b\ /\ 2} first, then adds \texttt{a}.

\subsubsection{Question 5}\label{question-5-1}

Which import is required before using \texttt{sqrt(x)} in the textbook
style?

Answer

\begin{Shaded}
\begin{Highlighting}[]
\ImportTok{from}\NormalTok{ math }\ImportTok{import}\NormalTok{ sqrt}
\end{Highlighting}
\end{Shaded}

\begin{center}\rule{0.5\linewidth}{0.5pt}\end{center}

\section{2.3 Problem Solving: First Do It By
Hand}\label{problem-solving-first-do-it-by-hand}

Before translating a problem into Python, work out at least one concrete
example manually.

The principle is:

\begin{quote}
If you cannot compute a solution yourself, it is unlikely that you can
write a reliable program that automates the computation.
\end{quote}

\begin{center}\rule{0.5\linewidth}{0.5pt}\end{center}

\subsection{Example -- Alternating Tiles Along a
Wall}\label{example-alternating-tiles-along-a-wall}

Suppose black and white tiles are placed along a wall with the
requirements:

\begin{itemize}
\tightlist
\item
  the first tile is black;
\item
  colors alternate;
\item
  the last tile must also be black.
\end{itemize}

Assume:

\begin{Shaded}
\begin{Highlighting}[]
\NormalTok{Total width = 100 inches}
\NormalTok{Tile width  =   5 inches}
\end{Highlighting}
\end{Shaded}

If we use 20 tiles, the last tile would be white. Instead, think of the
pattern as:

\begin{Shaded}
\begin{Highlighting}[]
\NormalTok{B | WB | WB | WB | ... | WB}
\end{Highlighting}
\end{Shaded}

There is one initial black tile, followed by black/white pairs.

The first tile uses \texttt{5} inches, leaving \texttt{95} inches.

Each pair uses:

\begin{Shaded}
\begin{Highlighting}[]
\NormalTok{2 × 5 = 10 inches}
\end{Highlighting}
\end{Shaded}

Number of complete pairs:

\begin{Shaded}
\begin{Highlighting}[]
\NormalTok{95 // 10 = 9}
\end{Highlighting}
\end{Shaded}

Number of tiles:

\begin{Shaded}
\begin{Highlighting}[]
\NormalTok{1 + 2 × 9 = 19}
\end{Highlighting}
\end{Shaded}

The tiles occupy:

\begin{Shaded}
\begin{Highlighting}[]
\NormalTok{19 × 5 = 95 inches}
\end{Highlighting}
\end{Shaded}

Total unused space:

\begin{Shaded}
\begin{Highlighting}[]
\NormalTok{100 {-} 95 = 5 inches}
\end{Highlighting}
\end{Shaded}

Gap at each end:

\begin{Shaded}
\begin{Highlighting}[]
\NormalTok{5 / 2 = 2.5 inches}
\end{Highlighting}
\end{Shaded}

\subsubsection{General Algorithm}\label{general-algorithm}

\begin{Shaded}
\begin{Highlighting}[]
\NormalTok{number of pairs = integer part of}
\NormalTok{                  (total width {-} tile width) / (2 × tile width)}

\NormalTok{number of tiles = 1 + 2 × number of pairs}

\NormalTok{gap at each end =}
\NormalTok{    (total width {-} number of tiles × tile width) / 2}
\end{Highlighting}
\end{Shaded}

Doing the arithmetic by hand exposes the structure needed for the
algorithm.

\begin{center}\rule{0.5\linewidth}{0.5pt}\end{center}

\section{Worked Example 2.1 -- Computing Travel
Time}\label{worked-example-2.1-computing-travel-time}

A robot travels toward an item. It can move faster on a road than across
rocky terrain.

Inputs:

\begin{itemize}
\tightlist
\item
  horizontal distance to the item;
\item
  vertical distance to the item;
\item
  speed on the road;
\item
  speed on rocky terrain;
\item
  length of the first road segment.
\end{itemize}

The total travel time consists of two parts.

\subsubsection{Segment 1}\label{segment-1}

\begin{Shaded}
\begin{Highlighting}[]
\NormalTok{time1 = length1 / speed1}
\end{Highlighting}
\end{Shaded}

\subsubsection{Segment 2}\label{segment-2}

The second segment is the hypotenuse of a right triangle.

\[
\text{segment2Length}
=
\sqrt{(\text{xDistance}-\text{segment1Length})^2 + \text{yDistance}^2}
\]

Then:

\begin{Shaded}
\begin{Highlighting}[]
\NormalTok{time2 = segment2Length / speed2}
\end{Highlighting}
\end{Shaded}

Total:

\begin{Shaded}
\begin{Highlighting}[]
\NormalTok{totalTime = time1 + time2}
\end{Highlighting}
\end{Shaded}

Python version:

\begin{Shaded}
\begin{Highlighting}[]
\ImportTok{from}\NormalTok{ math }\ImportTok{import}\NormalTok{ sqrt}

\NormalTok{segment1Time }\OperatorTok{=}\NormalTok{ segment1Length }\OperatorTok{/}\NormalTok{ segment1Speed}
\NormalTok{segment2Length }\OperatorTok{=}\NormalTok{ sqrt((xDistance }\OperatorTok{{-}}\NormalTok{ segment1Length) }\OperatorTok{**} \DecValTok{2} \OperatorTok{+}\NormalTok{ yDistance }\OperatorTok{**} \DecValTok{2}\NormalTok{)}
\NormalTok{segment2Time }\OperatorTok{=}\NormalTok{ segment2Length }\OperatorTok{/}\NormalTok{ segment2Speed}
\NormalTok{totalTime }\OperatorTok{=}\NormalTok{ segment1Time }\OperatorTok{+}\NormalTok{ segment2Time}
\end{Highlighting}
\end{Shaded}

The worked example also reinforces the use of descriptive variable names
in the final program.

\begin{center}\rule{0.5\linewidth}{0.5pt}\end{center}

\subsection{Self Check 2.3}\label{self-check-2.3}

\subsubsection{Question 1}\label{question-1-2}

Why should a programmer work out a concrete example by hand before
coding?

Answer

Because a hand calculation helps reveal the required steps and exposes
logical errors before Python syntax is introduced.

\subsubsection{Question 2}\label{question-2-2}

In the tile example, why is floor division useful?

Answer

Because only complete tile pairs can be placed. A fractional pair is not
meaningful.

\subsubsection{Question 3}\label{question-3-2}

Which order is preferable?

A. Write code → guess the algorithm → test\\
B. Work out an example → derive the algorithm → write code\\
C. Write code → choose variable names → understand the problem

Answer

\textbf{B.}

\begin{center}\rule{0.5\linewidth}{0.5pt}\end{center}

\section{2.4 Strings}\label{strings}

Many programs process text rather than numbers.

A \textbf{string} is a sequence of characters.

Examples:

\begin{Shaded}
\begin{Highlighting}[]
\CommentTok{"Hello"}
\CommentTok{"Python"}
\CommentTok{"123 Main Street"}
\end{Highlighting}
\end{Shaded}

Characters include:

\begin{itemize}
\tightlist
\item
  letters;
\item
  digits;
\item
  punctuation;
\item
  spaces;
\item
  Unicode symbols.
\end{itemize}

\begin{center}\rule{0.5\linewidth}{0.5pt}\end{center}

\subsection{2.4.1 The String Type}\label{the-string-type}

A string can be stored in a variable.

\begin{Shaded}
\begin{Highlighting}[]
\NormalTok{greeting }\OperatorTok{=} \StringTok{"Hello"}
\BuiltInTok{print}\NormalTok{(greeting)}
\end{Highlighting}
\end{Shaded}

A \textbf{string literal} is a specific string written directly in
source code.

Python supports both single and double quotation marks:

\begin{Shaded}
\begin{Highlighting}[]
\CommentTok{"This is a string."}
\CommentTok{\textquotesingle{}So is this.\textquotesingle{}}
\end{Highlighting}
\end{Shaded}

This allows one quotation style to appear inside the other:

\begin{Shaded}
\begin{Highlighting}[]
\NormalTok{message }\OperatorTok{=} \StringTok{\textquotesingle{}He said "Hello"\textquotesingle{}}
\end{Highlighting}
\end{Shaded}

\subsubsection{String Length}\label{string-length}

Use \texttt{len()} to obtain the number of characters.

\begin{Shaded}
\begin{Highlighting}[]
\NormalTok{length }\OperatorTok{=} \BuiltInTok{len}\NormalTok{(}\StringTok{"World!"}\NormalTok{)}
\end{Highlighting}
\end{Shaded}

\texttt{length} becomes \texttt{6}.

Spaces count as characters.

The \textbf{empty string} contains zero characters:

\begin{Shaded}
\begin{Highlighting}[]
\CommentTok{""}
\CommentTok{\textquotesingle{}\textquotesingle{}}
\end{Highlighting}
\end{Shaded}

\begin{center}\rule{0.5\linewidth}{0.5pt}\end{center}

\subsection{2.4.2 Concatenation and
Repetition}\label{concatenation-and-repetition}

\subsubsection{\texorpdfstring{Concatenation
\texttt{+}}{Concatenation +}}\label{concatenation}

The \texttt{+} operator joins strings.

\begin{Shaded}
\begin{Highlighting}[]
\NormalTok{firstName }\OperatorTok{=} \StringTok{"Harry"}
\NormalTok{lastName }\OperatorTok{=} \StringTok{"Morgan"}
\NormalTok{name }\OperatorTok{=}\NormalTok{ firstName }\OperatorTok{+} \StringTok{" "} \OperatorTok{+}\NormalTok{ lastName}
\end{Highlighting}
\end{Shaded}

Result:

\begin{Shaded}
\begin{Highlighting}[]
\NormalTok{Harry Morgan}
\end{Highlighting}
\end{Shaded}

Both operands must be strings.

Incorrect:

\begin{Shaded}
\begin{Highlighting}[]
\CommentTok{"Agent "} \OperatorTok{+} \DecValTok{1729}
\end{Highlighting}
\end{Shaded}

\subsubsection{\texorpdfstring{Repetition
\texttt{*}}{Repetition *}}\label{repetition}

A string can be repeated by an integer factor.

\begin{Shaded}
\begin{Highlighting}[]
\NormalTok{dashes }\OperatorTok{=} \StringTok{"{-}"} \OperatorTok{*} \DecValTok{50}
\end{Highlighting}
\end{Shaded}

Another example:

\begin{Shaded}
\begin{Highlighting}[]
\NormalTok{message }\OperatorTok{=} \StringTok{"Echo. "}
\BuiltInTok{print}\NormalTok{(message }\OperatorTok{*} \DecValTok{5}\NormalTok{)}
\end{Highlighting}
\end{Shaded}

\begin{center}\rule{0.5\linewidth}{0.5pt}\end{center}

\subsection{2.4.3 Converting Between Numbers and
Strings}\label{converting-between-numbers-and-strings}

\subsubsection{Number → String}\label{number-string}

Use \texttt{str()}.

\begin{Shaded}
\begin{Highlighting}[]
\BuiltInTok{id} \OperatorTok{=} \DecValTok{1729}
\NormalTok{name }\OperatorTok{=} \StringTok{"Agent "} \OperatorTok{+} \BuiltInTok{str}\NormalTok{(}\BuiltInTok{id}\NormalTok{)}
\end{Highlighting}
\end{Shaded}

Result:

\begin{Shaded}
\begin{Highlighting}[]
\NormalTok{Agent 1729}
\end{Highlighting}
\end{Shaded}

\subsubsection{String → Integer}\label{string-integer}

Use \texttt{int()}.

\begin{Shaded}
\begin{Highlighting}[]
\BuiltInTok{id} \OperatorTok{=} \BuiltInTok{int}\NormalTok{(}\StringTok{"1729"}\NormalTok{)}
\end{Highlighting}
\end{Shaded}

\subsubsection{String → Floating-Point
Number}\label{string-floating-point-number}

Use \texttt{float()}.

\begin{Shaded}
\begin{Highlighting}[]
\NormalTok{price }\OperatorTok{=} \BuiltInTok{float}\NormalTok{(}\StringTok{"17.29"}\NormalTok{)}
\end{Highlighting}
\end{Shaded}

The string must represent a valid numeric literal for the requested
conversion.

\begin{center}\rule{0.5\linewidth}{0.5pt}\end{center}

\subsection{2.4.4 Strings and Characters}\label{strings-and-characters}

A string is a sequence of characters. Each character has an
\textbf{index}.

Python starts counting at \texttt{0}.

\begin{Shaded}
\begin{Highlighting}[]
\NormalTok{String:  H  a  r  r  y}
\NormalTok{Index:   0  1  2  3  4}
\end{Highlighting}
\end{Shaded}

Given:

\begin{Shaded}
\begin{Highlighting}[]
\NormalTok{name }\OperatorTok{=} \StringTok{"Harry"}
\end{Highlighting}
\end{Shaded}

we can access individual characters:

\begin{Shaded}
\begin{Highlighting}[]
\NormalTok{first }\OperatorTok{=}\NormalTok{ name[}\DecValTok{0}\NormalTok{]}
\NormalTok{last }\OperatorTok{=}\NormalTok{ name[}\DecValTok{4}\NormalTok{]}
\end{Highlighting}
\end{Shaded}

Results:

\begin{Shaded}
\begin{Highlighting}[]
\NormalTok{first = "H"}
\NormalTok{last  = "y"}
\end{Highlighting}
\end{Shaded}

The last valid index is:

\begin{Shaded}
\begin{Highlighting}[]
\BuiltInTok{len}\NormalTok{(name) }\OperatorTok{{-}} \DecValTok{1}
\end{Highlighting}
\end{Shaded}

A general way to get the last character is:

\begin{Shaded}
\begin{Highlighting}[]
\NormalTok{last }\OperatorTok{=}\NormalTok{ name[}\BuiltInTok{len}\NormalTok{(name) }\OperatorTok{{-}} \DecValTok{1}\NormalTok{]}
\end{Highlighting}
\end{Shaded}

Using an index outside the valid range causes an exception at run time.

\subsubsection{Example -- Building
Initials}\label{example-building-initials}

\begin{Shaded}
\begin{Highlighting}[]
\NormalTok{first }\OperatorTok{=} \StringTok{"Rodolfo"}
\NormalTok{second }\OperatorTok{=} \StringTok{"Sally"}
\NormalTok{initials }\OperatorTok{=}\NormalTok{ first[}\DecValTok{0}\NormalTok{] }\OperatorTok{+} \StringTok{"\&"} \OperatorTok{+}\NormalTok{ second[}\DecValTok{0}\NormalTok{]}
\BuiltInTok{print}\NormalTok{(initials)}
\end{Highlighting}
\end{Shaded}

Output:

\begin{Shaded}
\begin{Highlighting}[]
\NormalTok{R\&S}
\end{Highlighting}
\end{Shaded}

\begin{center}\rule{0.5\linewidth}{0.5pt}\end{center}

\subsection{2.4.5 String Methods}\label{string-methods}

A \textbf{method} is an operation associated with an object.

For a string object:

\begin{Shaded}
\begin{Highlighting}[]
\NormalTok{name }\OperatorTok{=} \StringTok{"John Smith"}
\end{Highlighting}
\end{Shaded}

we can call:

\begin{Shaded}
\begin{Highlighting}[]
\NormalTok{uppercaseName }\OperatorTok{=}\NormalTok{ name.upper()}
\NormalTok{lowercaseName }\OperatorTok{=}\NormalTok{ name.lower()}
\end{Highlighting}
\end{Shaded}

Useful methods:

{\def\LTcaptype{none} % do not increment counter
\begin{longtable}[]{@{}ll@{}}
\toprule\noalign{}
Method & Result \\
\midrule\noalign{}
\endhead
\bottomrule\noalign{}
\endlastfoot
\texttt{s.lower()} & Lowercase version of \texttt{s} \\
\texttt{s.upper()} & Uppercase version of \texttt{s} \\
\texttt{s.replace(old,\ new)} & New string with replacements \\
\end{longtable}
}

Example:

\begin{Shaded}
\begin{Highlighting}[]
\NormalTok{name2 }\OperatorTok{=}\NormalTok{ name.replace(}\StringTok{"John"}\NormalTok{, }\StringTok{"Jane"}\NormalTok{)}
\end{Highlighting}
\end{Shaded}

Result:

\begin{Shaded}
\begin{Highlighting}[]
\NormalTok{Jane Smith}
\end{Highlighting}
\end{Shaded}

These methods return new strings. They do not modify the original string
object.

\begin{center}\rule{0.5\linewidth}{0.5pt}\end{center}

\section{Special Topic 2.4 -- Character
Values}\label{special-topic-2.4-character-values}

Characters are represented internally by integer codes.

Use \texttt{ord()} to obtain the code of a character:

\begin{Shaded}
\begin{Highlighting}[]
\BuiltInTok{ord}\NormalTok{(}\StringTok{"H"}\NormalTok{)}
\end{Highlighting}
\end{Shaded}

Use \texttt{chr()} to convert a numeric code back to a character:

\begin{Shaded}
\begin{Highlighting}[]
\BuiltInTok{chr}\NormalTok{(}\DecValTok{97}\NormalTok{)}
\end{Highlighting}
\end{Shaded}

Examples:

\begin{Shaded}
\begin{Highlighting}[]
\BuiltInTok{print}\NormalTok{(}\StringTok{"The letter H has code"}\NormalTok{, }\BuiltInTok{ord}\NormalTok{(}\StringTok{"H"}\NormalTok{))}
\BuiltInTok{print}\NormalTok{(}\StringTok{"Code 97 represents"}\NormalTok{, }\BuiltInTok{chr}\NormalTok{(}\DecValTok{97}\NormalTok{))}
\end{Highlighting}
\end{Shaded}

\begin{center}\rule{0.5\linewidth}{0.5pt}\end{center}

\section{Special Topic 2.5 -- Escape
Sequences}\label{special-topic-2.5-escape-sequences}

An \textbf{escape sequence} begins with a backslash and represents a
special character inside a string.

Examples:

\begin{Shaded}
\begin{Highlighting}[]
\CommentTok{"You\textquotesingle{}re }\CharTok{\textbackslash{}"}\CommentTok{Welcome}\CharTok{\textbackslash{}"}\CommentTok{"}
\end{Highlighting}
\end{Shaded}

\begin{Shaded}
\begin{Highlighting}[]
\CommentTok{"C:}\CharTok{\textbackslash{}\textbackslash{}}\CommentTok{Temp}\CharTok{\textbackslash{}\textbackslash{}}\CommentTok{Secret.txt"}
\end{Highlighting}
\end{Shaded}

\begin{Shaded}
\begin{Highlighting}[]
\BuiltInTok{print}\NormalTok{(}\StringTok{"*}\CharTok{\textbackslash{}n}\StringTok{**}\CharTok{\textbackslash{}n}\StringTok{***"}\NormalTok{)}
\end{Highlighting}
\end{Shaded}

Output:

\begin{Shaded}
\begin{Highlighting}[]
\NormalTok{*}
\NormalTok{**}
\NormalTok{***}
\end{Highlighting}
\end{Shaded}

Common sequences:

{\def\LTcaptype{none} % do not increment counter
\begin{longtable}[]{@{}ll@{}}
\toprule\noalign{}
Escape sequence & Meaning \\
\midrule\noalign{}
\endhead
\bottomrule\noalign{}
\endlastfoot
\texttt{\textbackslash{}"} & Double quotation mark inside a
double-quoted string \\
\texttt{\textbackslash{}\textbackslash{}} & Backslash \\
\texttt{\textbackslash{}n} & Newline \\
\end{longtable}
}

\begin{center}\rule{0.5\linewidth}{0.5pt}\end{center}

\section{Computing \& Society 2.1 -- International Alphabets and
Unicode}\label{computing-society-2.1-international-alphabets-and-unicode}

Text processing is not limited to the English alphabet.

Different writing systems include:

\begin{itemize}
\tightlist
\item
  accented European alphabets;
\item
  Greek;
\item
  Cyrillic;
\item
  Hebrew;
\item
  Arabic;
\item
  Chinese characters;
\item
  Japanese and Korean writing systems;
\item
  symbols and emoji.
\end{itemize}

\textbf{Unicode} provides a common encoding system capable of
representing characters from writing systems around the world.

Python 3 strings support Unicode.

Conceptually:

\begin{Shaded}
\begin{Highlighting}[]
\NormalTok{text }\OperatorTok{=} \StringTok{"£100"}
\BuiltInTok{print}\NormalTok{(text[}\DecValTok{0}\NormalTok{])}
\end{Highlighting}
\end{Shaded}

Each string position corresponds to a Unicode character, not simply an
ASCII byte.

\begin{center}\rule{0.5\linewidth}{0.5pt}\end{center}

\subsection{Self Check 2.4}\label{self-check-2.4}

\subsubsection{Question 1}\label{question-1-3}

What is the value of:

\begin{Shaded}
\begin{Highlighting}[]
\BuiltInTok{len}\NormalTok{(}\StringTok{"Hello"}\NormalTok{)}
\end{Highlighting}
\end{Shaded}

Answer

\texttt{5}

\subsubsection{Question 2}\label{question-2-3}

What is the value of:

\begin{Shaded}
\begin{Highlighting}[]
\CommentTok{"Py"} \OperatorTok{+} \StringTok{"thon"}
\end{Highlighting}
\end{Shaded}

Answer

\begin{Shaded}
\begin{Highlighting}[]
\NormalTok{Python}
\end{Highlighting}
\end{Shaded}

\subsubsection{Question 3}\label{question-3-3}

What does this produce?

\begin{Shaded}
\begin{Highlighting}[]
\CommentTok{"{-}"} \OperatorTok{*} \DecValTok{5}
\end{Highlighting}
\end{Shaded}

Answer

\begin{Shaded}
\begin{Highlighting}[]
\NormalTok{{-}{-}{-}{-}{-}}
\end{Highlighting}
\end{Shaded}

\subsubsection{Question 4}\label{question-4-2}

Given:

\begin{Shaded}
\begin{Highlighting}[]
\NormalTok{s }\OperatorTok{=} \StringTok{"Python"}
\end{Highlighting}
\end{Shaded}

what are \texttt{s{[}0{]}} and \texttt{s{[}len(s)\ -\ 1{]}}?

Answer

\begin{Shaded}
\begin{Highlighting}[]
\NormalTok{s[0]            → "P"}
\NormalTok{s[len(s) {-} 1]   → "n"}
\end{Highlighting}
\end{Shaded}

\subsubsection{Question 5}\label{question-5-2}

Why does this fail?

\begin{Shaded}
\begin{Highlighting}[]
\CommentTok{"Agent "} \OperatorTok{+} \DecValTok{1729}
\end{Highlighting}
\end{Shaded}

Answer

String concatenation requires strings on both sides. Convert the integer
first:

\begin{Shaded}
\begin{Highlighting}[]
\CommentTok{"Agent "} \OperatorTok{+} \BuiltInTok{str}\NormalTok{(}\DecValTok{1729}\NormalTok{)}
\end{Highlighting}
\end{Shaded}

\begin{center}\rule{0.5\linewidth}{0.5pt}\end{center}

\section{2.5 Input and Output}\label{input-and-output}

Most useful programs receive data from a user and produce results based
on those inputs.

A common structure is:

\begin{Shaded}
\begin{Highlighting}[]
\NormalTok{Input}
\NormalTok{  ↓}
\NormalTok{Process}
\NormalTok{  ↓}
\NormalTok{Output}
\end{Highlighting}
\end{Shaded}

\begin{center}\rule{0.5\linewidth}{0.5pt}\end{center}

\subsection{2.5.1 User Input}\label{user-input}

Use \texttt{input()} to read text from the keyboard.

\begin{Shaded}
\begin{Highlighting}[]
\NormalTok{first }\OperatorTok{=} \BuiltInTok{input}\NormalTok{(}\StringTok{"Enter your first name: "}\NormalTok{)}
\end{Highlighting}
\end{Shaded}

The string passed to \texttt{input()} is the \textbf{prompt}.

Conceptually, \texttt{input()} performs these steps:

\begin{Shaded}
\begin{Highlighting}[]
\NormalTok{Display prompt}
\NormalTok{     ↓}
\NormalTok{Wait for keyboard input}
\NormalTok{     ↓}
\NormalTok{User presses Enter}
\NormalTok{     ↓}
\NormalTok{Return entered characters as a string}
\end{Highlighting}
\end{Shaded}

Example:

\begin{Shaded}
\begin{Highlighting}[]
\NormalTok{first }\OperatorTok{=} \BuiltInTok{input}\NormalTok{(}\StringTok{"Enter your first name: "}\NormalTok{)}
\NormalTok{second }\OperatorTok{=} \BuiltInTok{input}\NormalTok{(}\StringTok{"Enter your significant other\textquotesingle{}s first name: "}\NormalTok{)}

\NormalTok{initials }\OperatorTok{=}\NormalTok{ first[}\DecValTok{0}\NormalTok{] }\OperatorTok{+} \StringTok{"\&"} \OperatorTok{+}\NormalTok{ second[}\DecValTok{0}\NormalTok{]}
\BuiltInTok{print}\NormalTok{(initials)}
\end{Highlighting}
\end{Shaded}

\begin{center}\rule{0.5\linewidth}{0.5pt}\end{center}

\subsection{2.5.2 Numerical Input}\label{numerical-input}

\texttt{input()} always returns a string.

To read an integer:

\begin{Shaded}
\begin{Highlighting}[]
\NormalTok{userInput }\OperatorTok{=} \BuiltInTok{input}\NormalTok{(}\StringTok{"Please enter the number of bottles: "}\NormalTok{)}
\NormalTok{bottles }\OperatorTok{=} \BuiltInTok{int}\NormalTok{(userInput)}
\end{Highlighting}
\end{Shaded}

To read a floating-point value:

\begin{Shaded}
\begin{Highlighting}[]
\NormalTok{userInput }\OperatorTok{=} \BuiltInTok{input}\NormalTok{(}\StringTok{"Enter price per bottle: "}\NormalTok{)}
\NormalTok{price }\OperatorTok{=} \BuiltInTok{float}\NormalTok{(userInput)}
\end{Highlighting}
\end{Shaded}

These operations can also be combined:

\begin{Shaded}
\begin{Highlighting}[]
\NormalTok{bottles }\OperatorTok{=} \BuiltInTok{int}\NormalTok{(}\BuiltInTok{input}\NormalTok{(}\StringTok{"Please enter the number of bottles: "}\NormalTok{))}
\NormalTok{price }\OperatorTok{=} \BuiltInTok{float}\NormalTok{(}\BuiltInTok{input}\NormalTok{(}\StringTok{"Enter price per bottle: "}\NormalTok{))}
\end{Highlighting}
\end{Shaded}

\begin{center}\rule{0.5\linewidth}{0.5pt}\end{center}

\subsection{2.5.3 Formatted Output}\label{formatted-output}

The chapter uses Python's \texttt{\%} \textbf{string format operator} to
control how values are displayed.

\subsubsection{Floating-Point Values}\label{floating-point-values}

\begin{Shaded}
\begin{Highlighting}[]
\NormalTok{price }\OperatorTok{=} \FloatTok{1.215962441314554}
\BuiltInTok{print}\NormalTok{(}\StringTok{"}\SpecialCharTok{\%.2f}\StringTok{"} \OperatorTok{\%}\NormalTok{ price)}
\end{Highlighting}
\end{Shaded}

Output:

\begin{Shaded}
\begin{Highlighting}[]
\NormalTok{1.22}
\end{Highlighting}
\end{Shaded}

\subsubsection{Field Width}\label{field-width}

\begin{Shaded}
\begin{Highlighting}[]
\BuiltInTok{print}\NormalTok{(}\StringTok{"}\SpecialCharTok{\%10.2f}\StringTok{"} \OperatorTok{\%}\NormalTok{ price)}
\end{Highlighting}
\end{Shaded}

The value occupies a field of width \texttt{10} and is right-justified.

\subsubsection{Common Format Specifiers}\label{common-format-specifiers}

{\def\LTcaptype{none} % do not increment counter
\begin{longtable}[]{@{}ll@{}}
\toprule\noalign{}
Specifier & Meaning \\
\midrule\noalign{}
\endhead
\bottomrule\noalign{}
\endlastfoot
\texttt{\%d} & Integer \\
\texttt{\%f} & Floating-point value \\
\texttt{\%.2f} & Float with 2 digits after decimal point \\
\texttt{\%7.2f} & Width 7, 2 decimal places \\
\texttt{\%s} & String \\
\texttt{\%9s} & Right-justified string in width 9 \\
\texttt{\%-9s} & Left-justified string in width 9 \\
\texttt{\%\%} & Literal percent sign \\
\texttt{\%+5d} & Show sign for positive/negative integer \\
\end{longtable}
}

\subsubsection{Formatting Multiple
Values}\label{formatting-multiple-values}

\begin{Shaded}
\begin{Highlighting}[]
\NormalTok{quantity }\OperatorTok{=} \DecValTok{24}
\NormalTok{total }\OperatorTok{=} \FloatTok{17.29}

\BuiltInTok{print}\NormalTok{(}\StringTok{"Quantity: }\SpecialCharTok{\%d}\StringTok{ Total: }\SpecialCharTok{\%10.2f}\StringTok{"} \OperatorTok{\%}\NormalTok{ (quantity, total))}
\end{Highlighting}
\end{Shaded}

Values are matched to format specifiers from left to right.

\begin{center}\rule{0.5\linewidth}{0.5pt}\end{center}

\subsection{Example -- Price per Ounce}\label{example-price-per-ounce}

\begin{Shaded}
\begin{Highlighting}[]
\CommentTok{\#\#}
\CommentTok{\# This program prints the price per ounce for a six{-}pack of cans.}
\CommentTok{\#}

\NormalTok{CANS\_PER\_PACK }\OperatorTok{=} \DecValTok{6}

\NormalTok{packPrice }\OperatorTok{=} \BuiltInTok{float}\NormalTok{(}\BuiltInTok{input}\NormalTok{(}\StringTok{"Please enter the price for a six{-}pack: "}\NormalTok{))}
\NormalTok{canVolume }\OperatorTok{=} \BuiltInTok{float}\NormalTok{(}\BuiltInTok{input}\NormalTok{(}\StringTok{"Please enter the volume for each can (in ounces): "}\NormalTok{))}

\NormalTok{packVolume }\OperatorTok{=}\NormalTok{ canVolume }\OperatorTok{*}\NormalTok{ CANS\_PER\_PACK}
\NormalTok{pricePerOunce }\OperatorTok{=}\NormalTok{ packPrice }\OperatorTok{/}\NormalTok{ packVolume}

\BuiltInTok{print}\NormalTok{(}\StringTok{"Price per ounce: }\SpecialCharTok{\%8.2f}\StringTok{"} \OperatorTok{\%}\NormalTok{ pricePerOunce)}
\end{Highlighting}
\end{Shaded}

This program demonstrates the full \textbf{Input → Process → Output}
pattern.

\begin{center}\rule{0.5\linewidth}{0.5pt}\end{center}

\section{Programming Tip 2.4 -- Don't Wait to
Convert}\label{programming-tip-2.4-dont-wait-to-convert}

When reading numeric input, convert it immediately.

Less desirable:

\begin{Shaded}
\begin{Highlighting}[]
\NormalTok{unitPrice }\OperatorTok{=} \BuiltInTok{input}\NormalTok{(}\StringTok{"Enter the unit price: "}\NormalTok{)}
\NormalTok{price1 }\OperatorTok{=} \BuiltInTok{float}\NormalTok{(unitPrice)}
\NormalTok{price2 }\OperatorTok{=} \DecValTok{12} \OperatorTok{*} \BuiltInTok{float}\NormalTok{(unitPrice)}
\end{Highlighting}
\end{Shaded}

Better:

\begin{Shaded}
\begin{Highlighting}[]
\NormalTok{unitPriceInput }\OperatorTok{=} \BuiltInTok{input}\NormalTok{(}\StringTok{"Enter the unit price: "}\NormalTok{)}
\NormalTok{unitPrice }\OperatorTok{=} \BuiltInTok{float}\NormalTok{(unitPriceInput)}

\NormalTok{price1 }\OperatorTok{=}\NormalTok{ unitPrice}
\NormalTok{price2 }\OperatorTok{=} \DecValTok{12} \OperatorTok{*}\NormalTok{ unitPrice}
\end{Highlighting}
\end{Shaded}

Or combine the calls:

\begin{Shaded}
\begin{Highlighting}[]
\NormalTok{unitPrice }\OperatorTok{=} \BuiltInTok{float}\NormalTok{(}\BuiltInTok{input}\NormalTok{(}\StringTok{"Enter the unit price: "}\NormalTok{))}
\end{Highlighting}
\end{Shaded}

Immediate conversion reduces repetition and makes it less likely that a
string will accidentally be used in a numerical computation.

\begin{center}\rule{0.5\linewidth}{0.5pt}\end{center}

\section{HOW TO 2.1 -- Writing Simple
Programs}\label{how-to-2.1-writing-simple-programs}

The chapter provides a systematic method for converting a problem
statement into a Python program.

\subsection{Problem -- Vending Machine
Change}\label{problem-vending-machine-change}

A customer:

\begin{itemize}
\tightlist
\item
  selects an item;
\item
  inserts a bill;
\item
  receives the purchased item;
\item
  receives change in dollar coins and quarters.
\end{itemize}

Assume item prices are multiples of 25 cents.

\begin{center}\rule{0.5\linewidth}{0.5pt}\end{center}

\subsection{Step 1. Understand the Inputs and
Outputs}\label{step-1.-understand-the-inputs-and-outputs}

\subsubsection{Inputs}\label{inputs}

\begin{itemize}
\tightlist
\item
  bill denomination;
\item
  item price in pennies.
\end{itemize}

\subsubsection{Outputs}\label{outputs}

\begin{itemize}
\tightlist
\item
  number of dollar coins;
\item
  number of quarters.
\end{itemize}

\begin{center}\rule{0.5\linewidth}{0.5pt}\end{center}

\subsection{Step 2. Work Out an Example by
Hand}\label{step-2.-work-out-an-example-by-hand}

Suppose:

\begin{Shaded}
\begin{Highlighting}[]
\NormalTok{Bill inserted = $5.00}
\NormalTok{Item price    = $2.25}
\end{Highlighting}
\end{Shaded}

Change:

\begin{Shaded}
\begin{Highlighting}[]
\NormalTok{$5.00 {-} $2.25 = $2.75 = 275 pennies}
\end{Highlighting}
\end{Shaded}

So the machine returns:

\begin{Shaded}
\begin{Highlighting}[]
\NormalTok{2 dollar coins}
\NormalTok{3 quarters}
\end{Highlighting}
\end{Shaded}

\begin{center}\rule{0.5\linewidth}{0.5pt}\end{center}

\subsection{Step 3. Write Pseudocode}\label{step-3.-write-pseudocode}

\begin{Shaded}
\begin{Highlighting}[]
\NormalTok{change due = 100 × bill value {-} item price}

\NormalTok{dollar coins = change due // 100}

\NormalTok{change due = change due \% 100}

\NormalTok{quarters = change due // 25}
\end{Highlighting}
\end{Shaded}

\begin{center}\rule{0.5\linewidth}{0.5pt}\end{center}

\subsection{Step 4. Identify Variables and
Constants}\label{step-4.-identify-variables-and-constants}

Variables:

\begin{Shaded}
\begin{Highlighting}[]
\NormalTok{billValue}
\NormalTok{itemPrice}
\NormalTok{changeDue}
\NormalTok{dollarCoins}
\NormalTok{quarters}
\end{Highlighting}
\end{Shaded}

Constants:

\begin{Shaded}
\begin{Highlighting}[]
\NormalTok{PENNIES\_PER\_DOLLAR }\OperatorTok{=} \DecValTok{100}
\NormalTok{PENNIES\_PER\_QUARTER }\OperatorTok{=} \DecValTok{25}
\end{Highlighting}
\end{Shaded}

All values are integers because the computation uses floor division and
remainders.

\begin{center}\rule{0.5\linewidth}{0.5pt}\end{center}

\subsection{Step 5. Translate the Computation into
Python}\label{step-5.-translate-the-computation-into-python}

\begin{Shaded}
\begin{Highlighting}[]
\NormalTok{changeDue }\OperatorTok{=}\NormalTok{ PENNIES\_PER\_DOLLAR }\OperatorTok{*}\NormalTok{ billValue }\OperatorTok{{-}}\NormalTok{ itemPrice}

\NormalTok{dollarCoins }\OperatorTok{=}\NormalTok{ changeDue }\OperatorTok{//}\NormalTok{ PENNIES\_PER\_DOLLAR}
\NormalTok{changeDue }\OperatorTok{=}\NormalTok{ changeDue }\OperatorTok{\%}\NormalTok{ PENNIES\_PER\_DOLLAR}
\NormalTok{quarters }\OperatorTok{=}\NormalTok{ changeDue }\OperatorTok{//}\NormalTok{ PENNIES\_PER\_QUARTER}
\end{Highlighting}
\end{Shaded}

\begin{center}\rule{0.5\linewidth}{0.5pt}\end{center}

\subsection{Step 6. Add Input and
Output}\label{step-6.-add-input-and-output}

\begin{Shaded}
\begin{Highlighting}[]
\NormalTok{billValue }\OperatorTok{=} \BuiltInTok{int}\NormalTok{(}\BuiltInTok{input}\NormalTok{(}\StringTok{"Enter bill value: "}\NormalTok{))}
\NormalTok{itemPrice }\OperatorTok{=} \BuiltInTok{int}\NormalTok{(}\BuiltInTok{input}\NormalTok{(}\StringTok{"Enter item price in pennies: "}\NormalTok{))}

\BuiltInTok{print}\NormalTok{(}\StringTok{"Dollar coins: }\SpecialCharTok{\%6d}\StringTok{"} \OperatorTok{\%}\NormalTok{ dollarCoins)}
\BuiltInTok{print}\NormalTok{(}\StringTok{"Quarters:     }\SpecialCharTok{\%6d}\StringTok{"} \OperatorTok{\%}\NormalTok{ quarters)}
\end{Highlighting}
\end{Shaded}

\begin{center}\rule{0.5\linewidth}{0.5pt}\end{center}

\subsection{Step 7. Assemble the
Program}\label{step-7.-assemble-the-program}

\begin{Shaded}
\begin{Highlighting}[]
\CommentTok{\#\#}
\CommentTok{\# This program simulates a vending machine that gives change.}
\CommentTok{\#}

\NormalTok{PENNIES\_PER\_DOLLAR }\OperatorTok{=} \DecValTok{100}
\NormalTok{PENNIES\_PER\_QUARTER }\OperatorTok{=} \DecValTok{25}

\NormalTok{billValue }\OperatorTok{=} \BuiltInTok{int}\NormalTok{(}\BuiltInTok{input}\NormalTok{(}\StringTok{"Enter bill value: "}\NormalTok{))}
\NormalTok{itemPrice }\OperatorTok{=} \BuiltInTok{int}\NormalTok{(}\BuiltInTok{input}\NormalTok{(}\StringTok{"Enter item price in pennies: "}\NormalTok{))}

\NormalTok{changeDue }\OperatorTok{=}\NormalTok{ PENNIES\_PER\_DOLLAR }\OperatorTok{*}\NormalTok{ billValue }\OperatorTok{{-}}\NormalTok{ itemPrice}

\NormalTok{dollarCoins }\OperatorTok{=}\NormalTok{ changeDue }\OperatorTok{//}\NormalTok{ PENNIES\_PER\_DOLLAR}
\NormalTok{changeDue }\OperatorTok{=}\NormalTok{ changeDue }\OperatorTok{\%}\NormalTok{ PENNIES\_PER\_DOLLAR}
\NormalTok{quarters }\OperatorTok{=}\NormalTok{ changeDue }\OperatorTok{//}\NormalTok{ PENNIES\_PER\_QUARTER}

\BuiltInTok{print}\NormalTok{(}\StringTok{"Dollar coins: }\SpecialCharTok{\%6d}\StringTok{"} \OperatorTok{\%}\NormalTok{ dollarCoins)}
\BuiltInTok{print}\NormalTok{(}\StringTok{"Quarters:     }\SpecialCharTok{\%6d}\StringTok{"} \OperatorTok{\%}\NormalTok{ quarters)}
\end{Highlighting}
\end{Shaded}

The method is general:

\begin{Shaded}
\begin{Highlighting}[]
\NormalTok{Understand}
\NormalTok{→ Calculate by hand}
\NormalTok{→ Pseudocode}
\NormalTok{→ Variables/constants}
\NormalTok{→ Python computation}
\NormalTok{→ Input/output}
\NormalTok{→ Complete program}
\end{Highlighting}
\end{Shaded}

\begin{center}\rule{0.5\linewidth}{0.5pt}\end{center}

\section{Worked Example 2.2 -- Computing the Cost of
Stamps}\label{worked-example-2.2-computing-the-cost-of-stamps}

A stamp vending machine receives dollar bills and dispenses:

\begin{itemize}
\tightlist
\item
  first-class stamps;
\item
  penny stamps as change.
\end{itemize}

Assume the first-class stamp price is \texttt{49} cents, as in the
textbook example.

\subsubsection{Input}\label{input}

\begin{itemize}
\tightlist
\item
  number of dollars inserted.
\end{itemize}

\subsubsection{Outputs}\label{outputs-1}

\begin{itemize}
\tightlist
\item
  number of first-class stamps;
\item
  number of penny stamps.
\end{itemize}

\subsubsection{Hand Calculation for \$1}\label{hand-calculation-for-1}

\begin{Shaded}
\begin{Highlighting}[]
\NormalTok{100 // 49 = 2 first{-}class stamps}
\NormalTok{100 {-} 2 × 49 = 2 cents change}
\end{Highlighting}
\end{Shaded}

\subsubsection{General Computation}\label{general-computation}

\begin{Shaded}
\begin{Highlighting}[]
\NormalTok{FIRST\_CLASS\_STAMP\_PRICE }\OperatorTok{=} \DecValTok{49}

\NormalTok{firstClassStamps }\OperatorTok{=} \DecValTok{100} \OperatorTok{*}\NormalTok{ dollars }\OperatorTok{//}\NormalTok{ FIRST\_CLASS\_STAMP\_PRICE}
\NormalTok{change }\OperatorTok{=} \DecValTok{100} \OperatorTok{*}\NormalTok{ dollars }\OperatorTok{{-}}\NormalTok{ firstClassStamps }\OperatorTok{*}\NormalTok{ FIRST\_CLASS\_STAMP\_PRICE}
\end{Highlighting}
\end{Shaded}

Complete example:

\begin{Shaded}
\begin{Highlighting}[]
\NormalTok{FIRST\_CLASS\_STAMP\_PRICE }\OperatorTok{=} \DecValTok{49}

\NormalTok{dollars }\OperatorTok{=} \BuiltInTok{int}\NormalTok{(}\BuiltInTok{input}\NormalTok{(}\StringTok{"Enter number of dollars: "}\NormalTok{))}

\NormalTok{firstClassStamps }\OperatorTok{=} \DecValTok{100} \OperatorTok{*}\NormalTok{ dollars }\OperatorTok{//}\NormalTok{ FIRST\_CLASS\_STAMP\_PRICE}
\NormalTok{change }\OperatorTok{=} \DecValTok{100} \OperatorTok{*}\NormalTok{ dollars }\OperatorTok{{-}}\NormalTok{ firstClassStamps }\OperatorTok{*}\NormalTok{ FIRST\_CLASS\_STAMP\_PRICE}

\BuiltInTok{print}\NormalTok{(}\StringTok{"First class stamps: }\SpecialCharTok{\%6d}\StringTok{"} \OperatorTok{\%}\NormalTok{ firstClassStamps)}
\BuiltInTok{print}\NormalTok{(}\StringTok{"Penny stamps:       }\SpecialCharTok{\%6d}\StringTok{"} \OperatorTok{\%}\NormalTok{ change)}
\end{Highlighting}
\end{Shaded}

This example reinforces:

\begin{itemize}
\tightlist
\item
  named constants;
\item
  integer arithmetic;
\item
  floor division;
\item
  input conversion;
\item
  formatted output.
\end{itemize}

\begin{center}\rule{0.5\linewidth}{0.5pt}\end{center}

\section{Computing \& Society 2.2 -- Bugs in
Silicon}\label{computing-society-2.2-bugs-in-silicon}

Not all incorrect numerical results are caused by application
programmers.

The chapter discusses the historical \textbf{Pentium floating-point
division bug}. A defect in processor hardware caused a very small set of
floating-point calculations to produce incorrect results.

The broader lesson is important:

\begin{itemize}
\tightlist
\item
  software depends on hardware;
\item
  hardware itself can contain defects;
\item
  numerical computing relies on carefully designed standards and
  implementations;
\item
  testing can reveal problems at multiple levels of a computing system.
\end{itemize}

The chapter also notes later processor-security issues related to
speculative execution, reinforcing the idea that sophisticated
optimizations can introduce unexpected risks.

\begin{center}\rule{0.5\linewidth}{0.5pt}\end{center}

\subsection{Self Check 2.5}\label{self-check-2.5}

\subsubsection{Question 1}\label{question-1-4}

What type does \texttt{input()} return?

Answer

A string (\texttt{str}).

\subsubsection{Question 2}\label{question-2-4}

How can an integer be read directly from user input?

Answer

\begin{Shaded}
\begin{Highlighting}[]
\NormalTok{value }\OperatorTok{=} \BuiltInTok{int}\NormalTok{(}\BuiltInTok{input}\NormalTok{(}\StringTok{"Enter a value: "}\NormalTok{))}
\end{Highlighting}
\end{Shaded}

\subsubsection{Question 3}\label{question-3-4}

What does \texttt{\%.2f} mean in the chapter's formatting syntax?

Answer

Display a floating-point value with two digits after the decimal point.

\subsubsection{Question 4}\label{question-4-3}

Why should numerical input be converted immediately?

Answer

It avoids repeated conversions and reduces the chance of accidentally
using a string in arithmetic.

\subsubsection{Question 5}\label{question-5-3}

Which operators are particularly useful when computing change in whole
coin units?

Answer

Floor division \texttt{//} and remainder \texttt{\%}.

\begin{center}\rule{0.5\linewidth}{0.5pt}\end{center}

\section{TOOLBOX 2.1 -- Symbolic Processing with
SymPy}\label{toolbox-2.1-symbolic-processing-with-sympy}

This optional section illustrates the broader Python ecosystem.

\textbf{SymPy} is a package for symbolic mathematics. Instead of only
computing numeric values, it can manipulate mathematical expressions
symbolically.

Typical setup:

\begin{Shaded}
\begin{Highlighting}[]
\ImportTok{from}\NormalTok{ sympy }\ImportTok{import} \OperatorTok{*}
\end{Highlighting}
\end{Shaded}

Create a symbolic expression:

\begin{Shaded}
\begin{Highlighting}[]
\NormalTok{f }\OperatorTok{=}\NormalTok{ sympify(}\StringTok{"x ** 2 * sin(x)"}\NormalTok{)}
\end{Highlighting}
\end{Shaded}

SymPy can perform tasks such as:

\subsubsection{Expanding Expressions}\label{expanding-expressions}

\begin{Shaded}
\begin{Highlighting}[]
\NormalTok{expand((x }\OperatorTok{{-}} \DecValTok{1}\NormalTok{) }\OperatorTok{*}\NormalTok{ (x }\OperatorTok{+} \DecValTok{1}\NormalTok{))}
\end{Highlighting}
\end{Shaded}

\subsubsection{Solving Equations}\label{solving-equations}

\begin{Shaded}
\begin{Highlighting}[]
\NormalTok{solve(x }\OperatorTok{**} \DecValTok{2} \OperatorTok{+} \DecValTok{2} \OperatorTok{*}\NormalTok{ x }\OperatorTok{{-}} \DecValTok{8}\NormalTok{)}
\end{Highlighting}
\end{Shaded}

\subsubsection{Differentiation}\label{differentiation}

\begin{Shaded}
\begin{Highlighting}[]
\NormalTok{diff(f)}
\end{Highlighting}
\end{Shaded}

\subsubsection{Integration}\label{integration}

\begin{Shaded}
\begin{Highlighting}[]
\NormalTok{g }\OperatorTok{=}\NormalTok{ integrate(f)}
\end{Highlighting}
\end{Shaded}

\subsubsection{Substitution and Numerical
Evaluation}\label{substitution-and-numerical-evaluation}

\begin{Shaded}
\begin{Highlighting}[]
\NormalTok{result }\OperatorTok{=}\NormalTok{ g.subs(x, }\DecValTok{0}\NormalTok{).evalf()}
\end{Highlighting}
\end{Shaded}

\subsubsection{Plotting}\label{plotting}

\begin{Shaded}
\begin{Highlighting}[]
\NormalTok{plot(x }\OperatorTok{**} \DecValTok{2}\NormalTok{)}
\end{Highlighting}
\end{Shaded}

The main lesson is not to memorize all SymPy commands. It is to
recognize that Python packages provide specialized expertise that can be
reused inside Python programs.

\begin{center}\rule{0.5\linewidth}{0.5pt}\end{center}

\section{Chapter 2 Summary}\label{chapter-2-summary}

\subsection{Variables}\label{variables-1}

\begin{itemize}
\tightlist
\item
  A variable is a named storage location.
\item
  Assignment stores a value in a variable.
\item
  A variable is created when first assigned.
\item
  Reassignment replaces the previous value.
\item
  \texttt{=} is assignment, not mathematical equality.
\item
  Important numeric types include \texttt{int} and \texttt{float}.
\item
  Use descriptive names.
\item
  Use uppercase names for constants according to the chapter convention.
\item
  Use comments to explain intent.
\end{itemize}

\subsection{Arithmetic}\label{arithmetic-1}

Important operators:

\begin{Shaded}
\begin{Highlighting}[]
\OperatorTok{+}  \OperatorTok{{-}}  \OperatorTok{*}  \OperatorTok{/}  \OperatorTok{//}  \OperatorTok{\%}  \OperatorTok{**}
\end{Highlighting}
\end{Shaded}

Remember:

\begin{Shaded}
\begin{Highlighting}[]
\NormalTok{/   → normal division}
\NormalTok{//  → floor division}
\NormalTok{\%   → remainder}
\NormalTok{**  → power}
\end{Highlighting}
\end{Shaded}

Useful functions:

\begin{Shaded}
\begin{Highlighting}[]
\BuiltInTok{abs}\NormalTok{()}
\BuiltInTok{round}\NormalTok{()}
\BuiltInTok{min}\NormalTok{()}
\BuiltInTok{max}\NormalTok{()}
\end{Highlighting}
\end{Shaded}

The \texttt{math} module provides additional functions such as:

\begin{Shaded}
\begin{Highlighting}[]
\NormalTok{sqrt()}
\NormalTok{sin()}
\NormalTok{cos()}
\NormalTok{log()}
\end{Highlighting}
\end{Shaded}

\subsection{Problem Solving}\label{problem-solving}

Before coding:

\begin{Shaded}
\begin{Highlighting}[]
\NormalTok{Choose concrete values}
\NormalTok{→ solve by hand}
\NormalTok{→ identify the general pattern}
\NormalTok{→ write the algorithm}
\NormalTok{→ translate into Python}
\end{Highlighting}
\end{Shaded}

\subsection{Strings}\label{strings-1}

\begin{itemize}
\tightlist
\item
  Strings are sequences of characters.
\item
  \texttt{len()} returns string length.
\item
  \texttt{+} concatenates strings.
\item
  \texttt{*} repeats a string.
\item
  \texttt{str()}, \texttt{int()}, and \texttt{float()} perform
  conversions.
\item
  String indices begin at \texttt{0}.
\item
  Methods such as \texttt{upper()}, \texttt{lower()}, and
  \texttt{replace()} return new strings.
\end{itemize}

\subsection{Input and Output}\label{input-and-output-1}

\begin{itemize}
\tightlist
\item
  \texttt{input()} returns a string.
\item
  Convert numeric input using \texttt{int()} or \texttt{float()}.
\item
  Convert as soon as practical after reading the input.
\item
  The chapter uses the \texttt{\%} string-format operator to control
  output layout.
\end{itemize}

\subsection{Graphics}\label{graphics}

\begin{itemize}
\tightlist
\item
  Create a graphics window.
\item
  Obtain the canvas.
\item
  Draw shapes and text on the canvas.
\item
  The origin is at the upper-left corner.
\item
  Color can be set by name or RGB components.
\end{itemize}

\begin{center}\rule{0.5\linewidth}{0.5pt}\end{center}

\section{Summary Quiz}\label{summary-quiz}

\subsection{Question 1}\label{question-1-5}

What does an assignment statement do?

A. Tests equality\\
B. Stores a value in a variable\\
C. Prints a value\\
D. Creates a comment

Answer

\textbf{B}

\subsection{Question 2}\label{question-2-5}

Which value has type \texttt{float}?

A. \texttt{6}\\
B. \texttt{0}\\
C. \texttt{6.0}\\
D. \texttt{-4}

Answer

\textbf{C}

\subsection{Question 3}\label{question-3-5}

What is the result of:

\begin{Shaded}
\begin{Highlighting}[]
\DecValTok{17} \OperatorTok{//} \DecValTok{5}
\end{Highlighting}
\end{Shaded}

A. \texttt{2}\\
B. \texttt{3}\\
C. \texttt{3.4}\\
D. \texttt{5}

Answer

\textbf{B}

\subsection{Question 4}\label{question-4-4}

What is the result of:

\begin{Shaded}
\begin{Highlighting}[]
\DecValTok{17} \OperatorTok{\%} \DecValTok{5}
\end{Highlighting}
\end{Shaded}

A. \texttt{2}\\
B. \texttt{3}\\
C. \texttt{5}\\
D. \texttt{17}

Answer

\textbf{A}

\subsection{Question 5}\label{question-5-4}

What does this expression produce?

\begin{Shaded}
\begin{Highlighting}[]
\CommentTok{"Py"} \OperatorTok{+} \StringTok{"thon"}
\end{Highlighting}
\end{Shaded}

Answer

\begin{Shaded}
\begin{Highlighting}[]
\NormalTok{Python}
\end{Highlighting}
\end{Shaded}

\subsection{Question 6}\label{question-6}

Given:

\begin{Shaded}
\begin{Highlighting}[]
\NormalTok{s }\OperatorTok{=} \StringTok{"Hello"}
\end{Highlighting}
\end{Shaded}

what is \texttt{s{[}1{]}}?

A. \texttt{"H"}\\
B. \texttt{"e"}\\
C. \texttt{"l"}\\
D. Error

Answer

\textbf{B. \texttt{"e"}}

\subsection{Question 7}\label{question-7}

What type does \texttt{input()} return?

A. Always \texttt{int}\\
B. Always \texttt{float}\\
C. \texttt{str}\\
D. Depends on what the user types

Answer

\textbf{C. \texttt{str}}

\subsection{Question 8}\label{question-8}

Which statement reads a floating-point value from the user?

A. \texttt{price\ =\ input()}\\
B. \texttt{price\ =\ float(input("Price:\ "))}\\
C. \texttt{price\ =\ str(input())}\\
D. \texttt{price\ =\ print(input())}

Answer

\textbf{B}

\subsection{Question 9}\label{question-9}

Which is the best reason to use a named constant?

A. It makes the program run faster in every case.\\
B. It makes an unexplained numeric value easier to understand and
change.\\
C. It prevents Python from changing the value.\\
D. It replaces comments entirely.

Answer

\textbf{B}

\subsection{Question 10}\label{question-10}

What should usually happen before translating a non-trivial computation
into Python?

A. Choose random syntax.\\
B. Work through at least one example by hand.\\
C. Install a graphics package.\\
D. Convert every value to a string.

Answer

\textbf{B}

\begin{center}\rule{0.5\linewidth}{0.5pt}\end{center}

\section{Review Exercises}\label{review-exercises}

\subsection{Exercise 1. Trace
Assignments}\label{exercise-1.-trace-assignments}

Determine the final value of \texttt{mystery}:

\begin{Shaded}
\begin{Highlighting}[]
\NormalTok{mystery }\OperatorTok{=} \DecValTok{1}
\NormalTok{mystery }\OperatorTok{=} \DecValTok{1} \OperatorTok{{-}} \DecValTok{2} \OperatorTok{*}\NormalTok{ mystery}
\NormalTok{mystery }\OperatorTok{=}\NormalTok{ mystery }\OperatorTok{+} \DecValTok{1}
\end{Highlighting}
\end{Shaded}

Answer

\begin{Shaded}
\begin{Highlighting}[]
\NormalTok{Initial: mystery = 1}
\NormalTok{Second statement: mystery = 1 {-} 2 × 1 = {-}1}
\NormalTok{Third statement: mystery = {-}1 + 1 = 0}
\end{Highlighting}
\end{Shaded}

Final value: \texttt{0}.

\begin{center}\rule{0.5\linewidth}{0.5pt}\end{center}

\subsection{Exercise 2. Translate Mathematics into
Python}\label{exercise-2.-translate-mathematics-into-python}

Write Python expressions for:

\begin{enumerate}
\def\labelenumi{\arabic{enumi}.}
\tightlist
\item
  The average of \texttt{a} and \texttt{b}.
\item
  \(a^2+b^2\).
\item
  \(\sqrt{x^2+y^2}\).
\item
  \(PV(1+r/100)^n\).
\end{enumerate}

Sample Answer

\begin{Shaded}
\begin{Highlighting}[]
\NormalTok{(a }\OperatorTok{+}\NormalTok{ b) }\OperatorTok{/} \DecValTok{2}

\NormalTok{a }\OperatorTok{**} \DecValTok{2} \OperatorTok{+}\NormalTok{ b }\OperatorTok{**} \DecValTok{2}

\NormalTok{sqrt(x }\OperatorTok{**} \DecValTok{2} \OperatorTok{+}\NormalTok{ y }\OperatorTok{**} \DecValTok{2}\NormalTok{)}

\NormalTok{PV }\OperatorTok{*}\NormalTok{ (}\DecValTok{1} \OperatorTok{+}\NormalTok{ r }\OperatorTok{/} \DecValTok{100}\NormalTok{) }\OperatorTok{**}\NormalTok{ n}
\end{Highlighting}
\end{Shaded}

\begin{center}\rule{0.5\linewidth}{0.5pt}\end{center}

\subsection{Exercise 3. Floor Division and
Remainder}\label{exercise-3.-floor-division-and-remainder}

Given:

\begin{Shaded}
\begin{Highlighting}[]
\NormalTok{n }\OperatorTok{=} \DecValTok{1729}
\end{Highlighting}
\end{Shaded}

predict:

\begin{Shaded}
\begin{Highlighting}[]
\NormalTok{n }\OperatorTok{\%} \DecValTok{10}
\NormalTok{n }\OperatorTok{//} \DecValTok{10}
\NormalTok{n }\OperatorTok{\%} \DecValTok{100}
\NormalTok{n }\OperatorTok{\%} \DecValTok{2}
\end{Highlighting}
\end{Shaded}

Answer

\begin{Shaded}
\begin{Highlighting}[]
\NormalTok{9}
\NormalTok{172}
\NormalTok{29}
\NormalTok{1}
\end{Highlighting}
\end{Shaded}

\begin{center}\rule{0.5\linewidth}{0.5pt}\end{center}

\subsection{Exercise 4. String
Expressions}\label{exercise-4.-string-expressions}

Given:

\begin{Shaded}
\begin{Highlighting}[]
\NormalTok{s }\OperatorTok{=} \StringTok{"Hello"}
\NormalTok{t }\OperatorTok{=} \StringTok{"World"}
\end{Highlighting}
\end{Shaded}

find the values of:

\begin{Shaded}
\begin{Highlighting}[]
\BuiltInTok{len}\NormalTok{(s) }\OperatorTok{+} \BuiltInTok{len}\NormalTok{(t)}
\NormalTok{s[}\DecValTok{1}\NormalTok{] }\OperatorTok{+}\NormalTok{ s[}\DecValTok{2}\NormalTok{]}
\NormalTok{s[}\BuiltInTok{len}\NormalTok{(s) }\OperatorTok{//} \DecValTok{2}\NormalTok{]}
\NormalTok{s }\OperatorTok{+}\NormalTok{ t}
\NormalTok{t }\OperatorTok{+}\NormalTok{ s}
\NormalTok{s }\OperatorTok{*} \DecValTok{2}
\end{Highlighting}
\end{Shaded}

Answer

\begin{Shaded}
\begin{Highlighting}[]
\NormalTok{10}
\NormalTok{el}
\NormalTok{l}
\NormalTok{HelloWorld}
\NormalTok{WorldHello}
\NormalTok{HelloHello}
\end{Highlighting}
\end{Shaded}

\begin{center}\rule{0.5\linewidth}{0.5pt}\end{center}

\subsection{Exercise 5. Explain the
Difference}\label{exercise-5.-explain-the-difference}

Explain the differences among:

\begin{Shaded}
\begin{Highlighting}[]
\NormalTok{2}
\NormalTok{2.0}
\NormalTok{\textquotesingle{}2\textquotesingle{}}
\NormalTok{"2"}
\NormalTok{"2.0"}
\end{Highlighting}
\end{Shaded}

Answer

\begin{itemize}
\tightlist
\item
  \texttt{2} is an integer.
\item
  \texttt{2.0} is a floating-point number.
\item
  \texttt{\textquotesingle{}2\textquotesingle{}} and \texttt{"2"} are
  strings containing one character.
\item
  \texttt{"2.0"} is a string containing three characters.
\end{itemize}

\begin{center}\rule{0.5\linewidth}{0.5pt}\end{center}

\subsection{Exercise 6. Roundoff
Error}\label{exercise-6.-roundoff-error}

Consider:

\begin{Shaded}
\begin{Highlighting}[]
\NormalTok{purchase }\OperatorTok{=} \FloatTok{19.93}
\NormalTok{payment }\OperatorTok{=} \FloatTok{20.00}
\NormalTok{change }\OperatorTok{=}\NormalTok{ payment }\OperatorTok{{-}}\NormalTok{ purchase}
\BuiltInTok{print}\NormalTok{(change)}
\end{Highlighting}
\end{Shaded}

Explain why a result such as \texttt{0.07000000000000028} may appear.

Answer

Many decimal fractions cannot be represented exactly in binary
floating-point form. The internal approximations can produce a tiny
numerical difference. Format or round the result when displaying money.

\begin{center}\rule{0.5\linewidth}{0.5pt}\end{center}

\section{Practice Exercises}\label{practice-exercises}

\subsection{Exercise 1. Paper
Dimensions}\label{exercise-1.-paper-dimensions}

A letter-size sheet is \texttt{8.5\ ×\ 11} inches. There are
\texttt{25.4} millimeters per inch.

Write a program that displays its dimensions in millimeters. Use named
constants and comments.

\begin{center}\rule{0.5\linewidth}{0.5pt}\end{center}

\subsection{Exercise 2. Powers of a
Number}\label{exercise-2.-powers-of-a-number}

Read a number and display:

\begin{itemize}
\tightlist
\item
  its square;
\item
  its cube;
\item
  its fourth power.
\end{itemize}

Use \texttt{**} only for the fourth power.

\begin{center}\rule{0.5\linewidth}{0.5pt}\end{center}

\subsection{Exercise 3. Two Integers}\label{exercise-3.-two-integers}

Read two integers and print:

\begin{itemize}
\tightlist
\item
  sum;
\item
  difference;
\item
  product;
\item
  average;
\item
  absolute difference;
\item
  maximum;
\item
  minimum.
\end{itemize}

\begin{center}\rule{0.5\linewidth}{0.5pt}\end{center}

\subsection{Exercise 4. Split a Five-Digit
Integer}\label{exercise-4.-split-a-five-digit-integer}

Read a five-digit positive integer and display its digits separately.

Example:

\begin{Shaded}
\begin{Highlighting}[]
\NormalTok{Input: 16384}
\NormalTok{Output: 1 6 3 8 4}
\end{Highlighting}
\end{Shaded}

Use floor division and remainder.

\begin{center}\rule{0.5\linewidth}{0.5pt}\end{center}

\subsection{Exercise 5. First, Middle, and Last
Character}\label{exercise-5.-first-middle-and-last-character}

Read a word and print:

\begin{itemize}
\tightlist
\item
  first character;
\item
  last character;
\item
  character in the middle.
\end{itemize}

For an even-length word, use the character immediately before the
middle.

\begin{center}\rule{0.5\linewidth}{0.5pt}\end{center}

\subsection{Exercise 6. Monogram}\label{exercise-6.-monogram}

Read a three-part name such as:

\begin{Shaded}
\begin{Highlighting}[]
\NormalTok{Harold James Morgan}
\end{Highlighting}
\end{Shaded}

and display:

\begin{Shaded}
\begin{Highlighting}[]
\NormalTok{HJM}
\end{Highlighting}
\end{Shaded}

First write pseudocode before attempting Python.

\begin{center}\rule{0.5\linewidth}{0.5pt}\end{center}

\subsection{Exercise 7. Vending
Machine}\label{exercise-7.-vending-machine}

Modify the vending-machine algorithm so that change is returned using:

\begin{itemize}
\tightlist
\item
  quarters;
\item
  dimes;
\item
  nickels.
\end{itemize}

Assume the item price is a multiple of \texttt{5} cents.

\begin{center}\rule{0.5\linewidth}{0.5pt}\end{center}

\subsection{Exercise 8. Temperature and Dew
Point}\label{exercise-8.-temperature-and-dew-point}

Using the formula and constants provided by your instructor or the
textbook exercise, write a program that reads temperature and relative
humidity and computes an approximate dew point.

Focus on:

\begin{itemize}
\tightlist
\item
  importing \texttt{log} from \texttt{math};
\item
  translating the mathematical formula carefully;
\item
  using parentheses correctly.
\end{itemize}

\begin{center}\rule{0.5\linewidth}{0.5pt}\end{center}

\section{Open Exercises}\label{open-exercises}

\subsection{Exercise 1}\label{exercise-1}

Give three examples of \textbf{magic numbers} that might appear in a
real business program. Replace each with a meaningful named constant.

\subsection{Exercise 2}\label{exercise-2}

Explain why \texttt{=} in Python should be read as ``assign'' rather
than ``equals''. Give an example that would be impossible as an
algebraic equation but valid as a Python assignment.

\subsection{Exercise 3}\label{exercise-3}

Choose an everyday quantity that naturally requires:

\begin{itemize}
\tightlist
\item
  floor division;
\item
  remainder.
\end{itemize}

Describe a program that uses both operators.

Examples include:

\begin{itemize}
\tightlist
\item
  converting seconds to minutes and seconds;
\item
  converting cents to dollars and cents;
\item
  packing items into complete boxes.
\end{itemize}

\subsection{Exercise 4}\label{exercise-4}

Choose a short string and show every character together with its index.

\begin{center}\rule{0.5\linewidth}{0.5pt}\end{center}

\section{Key Terms}\label{key-terms}

{\def\LTcaptype{none} % do not increment counter
\begin{longtable}[]{@{}
  >{\raggedright\arraybackslash}p{(\linewidth - 2\tabcolsep) * \real{0.5000}}
  >{\raggedright\arraybackslash}p{(\linewidth - 2\tabcolsep) * \real{0.5000}}@{}}
\toprule\noalign{}
\begin{minipage}[b]{\linewidth}\raggedright
Term
\end{minipage} & \begin{minipage}[b]{\linewidth}\raggedright
Meaning
\end{minipage} \\
\midrule\noalign{}
\endhead
\bottomrule\noalign{}
\endlastfoot
Variable & Named storage location \\
Assignment & Storing a value in a variable \\
Assignment operator & \texttt{=} \\
Initialization & Giving a variable its first value \\
Data type & Determines representation and valid operations \\
\texttt{int} & Integer type \\
\texttt{float} & Floating-point type \\
Literal & Value written directly in source code \\
Variable name & Identifier used to refer to a variable \\
Constant & Value intended to remain unchanged \\
Magic number & Unexplained numeric literal in code \\
Comment & Human-readable explanation ignored by Python \\
Operator & Symbol that performs an operation \\
Expression & Combination of values, variables, operators, and calls \\
Precedence & Rules controlling evaluation order \\
Exponentiation & Power operation, \texttt{**} \\
Floor division & Division with \texttt{//} \\
Modulus / remainder & \texttt{\%} operation \\
Function & Reusable operation that can receive arguments and return a
value \\
Standard library & Built-in collection of reusable modules \\
Module & Group of related functions and data definitions \\
Roundoff error & Small numerical error from finite floating-point
representation \\
String & Sequence of characters \\
String literal & String written directly in source code \\
Concatenation & Joining strings with \texttt{+} \\
Repetition & Repeating a string with \texttt{*} \\
Index & Position of a character in a string \\
Method & Operation associated with an object \\
Escape sequence & Backslash sequence representing a special character \\
Unicode & Character encoding standard for world writing systems \\
Prompt & Message that tells the user what input is expected \\
Input & Data supplied to a program \\
Output & Result produced by a program \\
Format specifier & Pattern controlling output representation \\
Canvas & Drawing area inside a graphics window \\
RGB & Red-green-blue color representation \\
Bounding box & Rectangle used to position an oval or other object \\
\end{longtable}
}

\begin{center}\rule{0.5\linewidth}{0.5pt}\end{center}

\section{Study Guidance}\label{study-guidance}

After Chapter 2, students should be able to write small programs that
manipulate numeric and text data. The most important habit is not
memorizing isolated syntax, but following a disciplined process:

\begin{Shaded}
\begin{Highlighting}[]
\NormalTok{Understand the task}
\NormalTok{→ choose meaningful variables and constants}
\NormalTok{→ work through an example by hand}
\NormalTok{→ write the computation clearly}
\NormalTok{→ obtain input}
\NormalTok{→ convert it to the correct type}
\NormalTok{→ process the data}
\NormalTok{→ format and display the output}
\NormalTok{→ test with simple values}
\NormalTok{→ inspect unexpected results}
\NormalTok{→ improve the program}
\end{Highlighting}
\end{Shaded}

A useful practice cycle is:

\begin{quote}
\textbf{Predict → Calculate by hand → Write → Run → Explain → Modify →
Practice}
\end{quote}

\end{document}
